\section{Proofs for the globally anchored connection}
\label{sec:supplementary-anchored-connection}

This section proves the results of
\cref{sec:anchored-physical-connection}.  The proof uses the long-delay
fading-space splitting already established in
\cref{prop:frozen-voltage-history-splitting}; it does not replace the RFDE by
a current-state system.

\begin{proof}[Proof of \cref{prop:anchor-central-identity-equilibria}]
Every retained fold coordinate and every member of its depth-two delay hull
tends to zero with the common diagonal wedge.  Hence all their collective
coordinates lie in \(|r|<r_{\rm loc}\) after decreasing that wedge.  There
\(g_{\rm anc}=1\), so \cref{eq:anchored-recovery-law} is literally
\cref{eq:shared-rfde-w}.  The voltage equation was not changed.  This proves
the full-history identity and transfers every local statement listed in the
proposition.

At a constant history all current-minus-delay differences in
\cref{eq:shared-rfde-v} vanish.  The two identities in
\cref{eq:outer-critical-transverse,eq:outer-critical-curve} say precisely
that its remaining voltage field vanishes at
\((V_{0,N}(r),w_{0,N}(r))\).  At \(r=\pm r_A\), the recovery field vanishes
because \(g_{\rm anc}(\pm r_A)=0\), proving
\cref{eq:anchor-equilibria}.  Finally,
\[
 \sigma_\pm'=(\pm r_A-\delta^2\nu)g_{\rm anc}'(\pm r_A).
\]
The first factor has the sign of \(\pm r_A\) and the two derivatives have
the signs in \cref{eq:anchor-multiplier}; both products are negative.
\end{proof}

\subsection{The characteristic Schur count}

At either anchor abbreviate
\begin{equation}
 J_\pm=\mathcal J^{\rm fv}_N(\pm r_A),\qquad
 A_\pm(\lambda)
 =J_\pm+\delta^2K\sum_{k=0}^mL_{k,N}(\eta)
       \{1-e^{-\lambda\theta_k/\delta}\}.           \label{eq:anchor-characteristic-voltage-pencil}
\end{equation}

\begin{lemma}[Uniform bordered characteristic count]
\label[lemma]{lem:anchor-characteristic-count}
Under the hypotheses of \cref{thm:anchor-indices-manifolds}, the full
characteristic equation at \(E_N^\pm\) has one simple root in
\(|\lambda|\le C\delta^2\), namely
\cref{eq:anchor-slow-roots}.  Outside that disk its open-right-half-plane
root count equals the frozen-resource voltage count: zero at \(E_N^+\) and
one at \(E_N^-\).  All constants and contour resolvent bounds are uniform in
\(N\); the count below uses one common half-disk boundary and does not
separate individual long-delay clusters.
\end{lemma}

\begin{proof}
For a modal ansatz \((u,\omega)e^{\lambda t}\), the characteristic equations
are
\begin{align}
 (\lambda\Id-A_\pm(\lambda))u&=-\omega\mathbf1,\notag\\
 \lambda\omega&=\delta^2\sigma_\pm'\pi_N^Tu.        \label{eq:anchor-full-characteristic-system}
\end{align}
On a fixed disk about zero that excludes the voltage spectrum, eliminate
\(u\).  The scalar Schur function is
\begin{equation}
 \Delta_\pm(\lambda)
 =\lambda+\delta^2\sigma_\pm'\pi_N^T
   (\lambda\Id-A_\pm(\lambda))^{-1}\mathbf1.        \label{eq:anchor-slow-Schur-function}
\end{equation}
The exact critical-curve identities
\begin{equation}
 J_\pm V_{0,N}'(\pm r_A)=w_{0,N}'(\pm r_A)\mathbf1,
 \qquad \pi_N^TV_{0,N}'(\pm r_A)=1                 \label{eq:anchor-critical-resolvent-identity}
\end{equation}
give
\begin{equation}
 \pi_N^T(-J_\pm)^{-1}\mathbf1
 =-\frac1{w_{0,N}'(\pm r_A)}.                       \label{eq:anchor-zero-resolvent-value}
\end{equation}
The fixed-radius voltage splitting and the common layer bound imply, for
\(|\lambda|\le C\delta^2\),
\begin{equation}
 \left\|(\lambda\Id-A_\pm(\lambda))^{-1}
          -(-J_\pm)^{-1}\right\|\le C\delta^2.     \label{eq:anchor-small-disk-resolvent}
\end{equation}
Indeed
\(|1-e^{-\lambda\theta_k/\delta}|\le C\delta\) on this disk, so the delayed
part is even \(O(\delta^3)\); the displayed weaker bound is convenient.
Thus
\begin{equation}
 \Delta_\pm(\lambda)
 =\lambda-\delta^2\frac{\sigma_\pm'}
                         {w_{0,N}'(\pm r_A)}
       +O(\delta^4),\qquad
 \Delta_\pm'(\lambda)=1+O(\delta^2).               \label{eq:anchor-Schur-expansion}
\end{equation}
Rouch\'{e}'s theorem on a circle of radius \(C\delta^2\), followed by the
implicit-function estimate, proves the unique slow root and
\cref{eq:anchor-slow-roots}.  The sign follows from
\cref{eq:outer-eigen-identity,eq:anchor-source-derivative-sign}.

It remains to prove the index without separating the long-delay
characteristic clusters.  We use a bordered homotopy.  For
\(0\le\varepsilon\le\delta^2\), replace the second row of
\cref{eq:anchor-full-characteristic-system} by
\(\lambda\omega=\varepsilon\sigma_\pm'\pi_N^Tu\) and set
\begin{equation}
 \Delta_{\pm,\varepsilon}(\lambda)
 =\lambda+\varepsilon\sigma_\pm'\pi_N^T
   (\lambda\Id-A_\pm(\lambda))^{-1}\mathbf1.       \label{eq:anchor-homotopy-Schur}
\end{equation}
Although this quotient is written only off the voltage roots, the product
\begin{equation}
 D_{\pm,\varepsilon}(\lambda)
 =\det(\lambda\Id-A_\pm(\lambda))
       \Delta_{\pm,\varepsilon}(\lambda)              \label{eq:anchor-bordered-characteristic-function}
\end{equation}
extends to the entire bordered characteristic determinant.  Thus poles in
the displayed Schur quotient create no missing or additional roots.

We record the boundary estimate used in the homotopy.  The fixed-anchor
current matrices have a common hyperbolicity margin.  The current-state
dichotomy in \cref{eq:frozen-current-attracting,eq:frozen-current-repelling}
and its Laplace transform give, for one \(C_A\) independent of
\(N,\delta,\eta\),
\begin{equation}
 \sup_{\xi\in\mathbb R}
 \|(i\xi\Id-A_\pm(i\xi))^{-1}\|\le C_A. \label{eq:anchor-imaginary-voltage-resolvent}
\end{equation}
Here is a direct verification that does not infer the estimate from a
spectral abscissa.  On \(|\xi|\le R_A\), the resolvent of
\(i\xi\Id-J_\pm\) is uniformly bounded by the fixed-anchor current
dichotomy, while
\(\|\delta^2K\sum_kL_{k,N}(1-e^{-i\xi\theta_k/\delta})\|\le
2b_*\delta^2\); a Neumann series applies after decreasing
\(\delta_0\).  Choose \(R_A\) larger than twice the common norms of
\(J_\pm\) and the delay layers.  For \(|\xi|>R_A\), factor
\(i\xi\) and use a second Neumann series.  This proves
\cref{eq:anchor-imaginary-voltage-resolvent} and, in particular, proves
that the voltage determinant has no imaginary-axis zero.

The same two-region argument shows that every zero of
\(D_{\pm,\varepsilon}\) in \(\Re\lambda\ge0\) lies in a
fixed half disk: for \(\Re\lambda\ge0\), all delay exponentials have
modulus at most one, and the current and scalar rows cannot balance a
\(\lambda\) of sufficiently large modulus.  Hence roots cannot escape
through infinity during the homotopy.

It remains only to exclude a crossing of the imaginary axis.  Uniformly
for the imaginary segment \(\lambda=i\xi\), \(|\xi|\le r_*\), and
\(0\le\varepsilon\le\delta^2\), the calculation leading to
\cref{eq:anchor-Schur-expansion} gives
\begin{equation}
 \Delta_{\pm,\varepsilon}(\lambda)
 =\lambda-\varepsilon
       \frac{\sigma_\pm'}{w_{0,N}'(\pm r_A)}
 +\varepsilon\mathcal R_\pm(\lambda,\delta),
 \qquad
 |\mathcal R_\pm(\lambda,\delta)|\le C|\lambda|,
 \qquad
 |\partial_\xi\mathcal R_\pm(i\xi,\delta)|\le C. \label{eq:anchor-homotopy-small-expansion}
\end{equation}
This estimate is deliberately restricted to the pole-free imaginary
segment supplied by \cref{eq:anchor-imaginary-voltage-resolvent}; no Schur
quotient is claimed analytic on a fixed complex disk containing long-delay
voltage roots.
The quotient \(\sigma_\pm'/w_{0,N}'(\pm r_A)\) is separated
from zero.  Consequently there is no imaginary zero in
\(|\xi|\le C_0\varepsilon\) when \(\varepsilon>0\): the real
part of \cref{eq:anchor-homotopy-small-expansion} has the fixed nonzero
sign of its second term.  For
\(C_0\varepsilon\le|\xi|\le r_*\), choose \(C_0>2C_A|\sigma_\pm'|\);
then \(|i\xi|\) dominates the second term in
\cref{eq:anchor-homotopy-Schur}.  On \(|\xi|\ge r_*\), first
decrease \(\delta_0\) and use
\cref{eq:anchor-imaginary-voltage-resolvent}.  These three regions
exclude every imaginary-axis zero for \(0<\varepsilon\le\delta^2\).

At \(\varepsilon=0\),
\(D_{\pm,0}(\lambda)=\lambda\det(\lambda\Id-A_\pm(\lambda))\).
For \(\varepsilon>0\), its zero at the origin moves to
\(\lambda=\varepsilon\sigma_\pm'/w_{0,N}'(\pm r_A)
+O(\varepsilon\delta^2+\varepsilon^2)\).  Remove a small half-disk around
this moving simple zero and apply the argument principle on the remaining
common right-half-disk boundary.  The small boundary counts exactly that
one root, while the outer boundary is uniformly root free.  The preceding
no-crossing estimate therefore
show that every other open-right-half-plane root count equals the voltage
count in \cref{prop:frozen-voltage-history-splitting}.  At the negative
anchor the real simple strong root satisfies, in original time,
\begin{equation}
 \lambda_{\mathrm{su},N}^{\mathrm{orig}}
 =\delta\lambda^u_{N,\delta,-r_A,\eta}+O(\delta^2)
 =\lambda_{c,N}(-r_A)+O(\delta^2),          \label{eq:anchor-strong-root-original-time}
\end{equation}
and remains real because it is the continuation of a real algebraically
simple voltage root.  This proves the asserted counts without a contour
exhaustion or an individual long-delay cluster separation.
\end{proof}

\begin{lemma}[Scaled local charts at the two anchors]
\label[lemma]{lem:anchor-scaled-local-manifolds}
Let \(\tau_\delta=\theta_m/\delta\) and equip
\(\mathcal H_{N,\delta}=C([-\tau_\delta,0];\mathbb R^N)\times\mathbb R\)
with the voltage fading norm obtained from
\cref{eq:frozen-history-norm} by the change \(s=\delta t\).  Put
\begin{equation}
 h_\pm=J_\pm^{-1}\mathbf1
       =\frac{V_{0,N}'(\pm r_A)}{w_{0,N}'(\pm r_A)},
 \qquad
 \|(\phi,\omega)\|_{\pm,\delta}^{\rm ph}
 =\|\phi-\omega h_\pm\|_{N,\pm r_A,\delta}^{\sharp}
       +|\omega|.                                  \label{eq:anchor-phase-norm}
\end{equation}
The linearized anchored semiflow has invariant projections
\(P_\pm^s,P_\pm^{\rm sl},P_\pm^{\rm su}\), with the last projection zero
at the positive anchor, such that
\begin{align}
 \|P_\pm^j\|&\le C,\notag\\
 \|T_+(t)P_+^s\|&\le Ce^{-c\delta\kappa_A t},
 &\|T_+(-t)P_+^{\rm sl}\|&\le Ce^{-c\delta^2t},\notag\\
 \|T_-(t)(P_-^s+P_-^{\rm sl})\|&\le Ce^{-c\delta^2t},
 &\|T_-(-t)P_-^{\rm su}\|&\le Ce^{-c t},          \label{eq:anchor-full-linear-trichotomy}
\end{align}
for \(t\ge0\), where
\(\kappa_A=\min_\pm\kappa_{N,\delta}(\pm r_A)\).
The projections and their first \((\nu,\eta)\)-derivatives have
norm at most \(C\delta^{-M_{\rm sp}}\) for one fixed
\(M_{\rm sp}\).  On the phase ball of radius
\(\rho_{A,\delta}=c_A\delta^4\), the Lyapunov--Perron maps for
\(W^u(E_N^+)\) and \(W^s(E_N^-)\), after the slow coordinate is divided
by \(\rho_{A,\delta}\), have a common contraction constant at most
\(1/3\).  Their first differentiated fixed-point equations have inverse
norm at most \(3/2\) and only a fixed polynomial loss in \(\delta^{-1}\).
All assertions are uniform in \(N\).  This local theorem does not select a
distinguished orbit on the stable sheet.  The chart-center point used below
is instead defined by the exact half-line construction, so no local stable
coordinate is continued backward through the RFDE.
\end{lemma}

\begin{proof}
At zero recovery feedback, meaning that the second row in
\cref{eq:anchor-full-characteristic-system} is replaced by
\(\dot\omega=0\), the phase isomorphism
\begin{equation}
 \mathcal C_{\pm,0}(\phi,\omega)
   =(\phi-\omega h_\pm,\omega)                         \label{eq:anchor-zero-feedback-phase-isomorphism}
\end{equation}
conjugates the variational semiflow to the direct product of the
frozen-resource voltage RFDE and one scalar center line.  This statement is
on the actual phase space \(C\times\mathbb R\): it subtracts the current
resource from the stored voltage history and does not introduce a fictitious
resource history.  It is exact because the resource is constant at zero
feedback, \(J_\pm h_\pm=\mathbf1\), and delayed differences annihilate the
constant history \(h_\pm\).

Turning on \(0\le\varepsilon\le\delta^2\) adds the bounded retarded
vector-field row
\((\phi,\omega)\mapsto(0,\varepsilon\sigma_\pm'\pi_N^T\phi(0))\).
Variation of constants in the norm
\cref{eq:anchor-phase-norm}, together with
\cref{eq:frozen-attracting-history-estimate,eq:frozen-repelling-stable-estimate,eq:frozen-unstable-inverse-estimate},
has perturbation number
\begin{equation}
 q_{\rm sp}
 \le C\delta^2
   \left(1+\frac{1}{\delta\kappa_A}\right)=o(1).  \label{eq:anchor-full-splitting-smallness}
\end{equation}
The center graph is obtained with a neutral boundary row rather than by
integrating the center kernel to infinity.  The graph equations for the
stable, voltage-unstable, and center columns are triangular; solving them
in that order gives bounded projections and a one-dimensional center
bundle.  The small-circle Schur calculation in
\cref{eq:anchor-Schur-expansion} identifies this bundle with the simple
slow eigenspace.  Its sign, and the half-plane count in
\cref{lem:anchor-characteristic-count}, convert the center bundle into the
splittings displayed in \cref{eq:anchor-full-linear-trichotomy}.  Notice
that the slow Green integral is \(O(\delta^{-2})\); this is recorded
rather than hidden in a claim of a common spectral gap.

For parameter jets, differentiate the three graph equations just described
once.
The layer map is affine in \(\eta\), its operator-sum norm is common, and
the anchor source and current matrices have common derivatives.  The first
derivative solves
\begin{equation}
 (I-D\mathcal T_{\rm sp})X_p
     =\partial_p\mathcal T_{\rm sp},                \label{eq:anchor-spectral-first-jet}
\end{equation}
Every Green convolution costs at most \(C\delta^{-2}\).  Thus one fixed
power \(M_{\rm sp}\) bounds all three projections and their first
rectangular jets.
This gives the claimed strong-to-weak bounds without differentiating a
moving evaluation atom or a family of long-delay contours.

In compatible solution-manifold coordinates centered at either equilibrium,
write the nonlinear equation as
\(x'=\mathcal A_\pm x+\mathcal N_\pm(x;\nu,\eta)\).  On a common
unit phase ball,
\begin{equation}
 \mathcal N_\pm(0)=D_x\mathcal N_\pm(0)=0,
 \qquad
 \|D_x\mathcal N_\pm(x)\|\le C\|x\|_{\pm,\delta}^{\rm ph},
 \qquad \|D_x^2\mathcal N_\pm(x)\|\le C.          \label{eq:anchor-local-nonlinear-bounds}
\end{equation}
The standard past Lyapunov--Perron formula on
\(P_+^{\rm sl}\oplus P_+^s\) and the future mixed formula on
\((P_-^s+P_-^{\rm sl})\oplus P_-^{\rm su}\) use exactly the kernels in
\cref{eq:anchor-full-linear-trichotomy}.  Their Lipschitz number on a ball
of radius \(\rho\) is at most
\(C\delta^{-2}\rho\).  With
\(\rho=c_A\delta^4\) and then \(c_A\) small it is at most \(1/3\).
The differentiated equation has the form
\begin{equation}
 (I-D\mathcal T_{\rm nl})x_p
   =\partial_p\mathcal T_{\rm nl}.                  \label{eq:anchor-local-manifold-jets}
\end{equation}
Its inverse norm is at most \(3/2\); the projection bounds and the
finite Green factors leave only a fixed polynomial loss.  This proves the
local graph assertions.

The first differentiated graph equation and the fixed parameter
trivializations give the stated \(C^1\) family of local sheets.  No
one-dimensional stable submanifold is selected here.  This proves the
lemma.
\end{proof}

\begin{proof}[Proof of \cref{thm:anchor-indices-manifolds}]
The spectral assertions are \cref{lem:anchor-characteristic-count}.  The
sign identity \cref{eq:outer-eigen-identity} gives
\(w_{0,N}'(r)<0\) at \(r=r_A\) and
\(w_{0,N}'(r)>0\) at \(r=-r_A\), after decreasing \(r_A\).  Together with
\(\sigma_\pm'<0\), this proves the signs of the slow roots.  The scaled
construction in \cref{lem:anchor-scaled-local-manifolds} gives a
one-dimensional \(W^u(E_N^+)\) and a codimension-one \(W^s(E_N^-)\) on
the compatible solution manifold.  For the fixed-system
invariant-manifold theorem used inside that uniform construction, see
\citet[Chs.~7 and~10]{hale1993introduction}.

The slow eigenvectors converge, in the fixed graph chart, to
\begin{equation}
 \left(V_{0,N}'(\pm r_A),w_{0,N}'(\pm r_A)\right),  \label{eq:anchor-slow-eigenvector}
\end{equation}
whose stationary current coordinate is one.  The normalized exact slow
eigenhistory therefore has stationary current coordinate
\(1+O(\delta^2)\).  The inward unstable hit at the scale
\cref{eq:anchor-local-manifold-scale} is transverse.  The negative local
sheet is represented in a fixed phase chart through the chart-center point;
no monotone \(r\)-clock is assigned to that sheet.  A point on the local unstable
manifold has the unique complete negative-time orbit supplied by the
unstable-manifold chart.  The stable manifold is defined by forward
convergence and makes no backward-completeness assertion.

Applying the event implicit-function theorem inside the scaled chart costs
only a fixed power of \(\delta^{-1}\), because the section speed is nonzero
there.  Equations
\cref{eq:anchor-spectral-first-jet,eq:anchor-local-manifold-jets}
give the first parameter derivatives and the asserted polynomial bound.  This
proves the stated rectangular regularity and uniformity.
\end{proof}

\subsection{Fold-time half-lines and the physical stable sheet}
\label{sec:anchored-fold-time-half-lines}

All constructions in this subsection use physical fold time \(s=\delta t\).
Put
\begin{equation}
 \mathscr X_N=C([-\theta_m,0];\mathbb R^N),
 \qquad
 \mathscr Y_N=\mathscr X_N\times\mathbb R.          \label{eq:anchor-fixed-delay-space}
\end{equation}
If \(\Phi\) is an original-time voltage history, its fold-time
representative is
\begin{equation}
 (\mathcal R_\delta\Phi)(\vartheta)
 =\Phi(\vartheta/\delta),\qquad -\theta_m\le\vartheta\le0.
                                                               \label{eq:anchor-history-rescaling}
\end{equation}
Thus the delayed values below are literally \(v(s-\theta_k)\).  The second
factor in \(\mathscr Y_N\) is the current resource value; no resource
history is adjoined.

For later reference, set
\[
 \mathscr V_N(v,w)
 =\frac23\mathbf1-w\mathbf1
  -c_N\circ(v-\mathbf1)^{\circ2}
  -\frac\beta3(v-\mathbf1)^{\circ3}
  +D(P_N-\Id)v .
\]
The raw fold-time equation is
\begin{align}
 v'(s)
 &=\delta^{-1}\mathscr V_N(v(s),w(s))
   +\delta K\sum_{k=0}^mL_{k,N}(\eta)
       \{v(s)-v(s-\theta_k)\},\notag\\
 w'(s)
 &=\delta\,\sigma_{\rm anc}
       \bigl(\pi_N^T(v(s)-\mathbf1);\delta,\nu\bigr).
                                                               \label{eq:anchor-raw-fold-time-rfde}
\end{align}
Denote the right-hand side, acting on a literal sliding history and the
current resource, by
\(\mathscr F^{\rm f}_{N,\delta,p}:\mathscr Y_N\to\mathbb R^{N+1}\).
For a current path write
\(Z_s=(\vartheta\mapsto v(s+\vartheta),w(s))\).
The weak half-line domain is
\begin{equation}
 \mathscr D_-^0
 =\left\{(v,w):
 \begin{array}{l}
 v\in C([-\theta_m,\infty);\mathbb R^N),\
 w\in C([0,\infty);\mathbb R),\\
 (v,w)\in W^{1,\infty}_{\rm loc}((0,\infty))
 \end{array}\right\}.                                \label{eq:anchor-half-line-weak-domain}
\end{equation}
Every occurrence of a delayed path in what follows means the literal
sliding history \(v_s\).

Fix \(0<r_{\rm seed}<r_{\rm in}<r_{\rm loc}\).  The preliminary negative
seed section is
\(\Sigma_-^{\rm seed}=\{r=-r_{\rm seed}\}\), while the strong entry
section is \(\Sigma_-^{\rm ent}=\{r=-r_{\rm in}\}\).  Their separation is
fixed.  An exact orbit starting at the seed section therefore accumulates
more than three complete delay buffers before it reaches the entry section
when \(\delta\) is small.  Fix also the corresponding positive inner
section at \(r=r_{\rm in}\).  All three sections stay a fixed distance from
the fold and the anchors.  Their physical event speeds on the declared
channels have a common nonzero margin after the channel estimates below.
Write \(\ell_{\rm seed}\) and \(\ell_{\rm ph}\) for their constant phase
rows, respectively.

\begin{lemma}[Exact fold-time anchor half-lines]
\label[lemma]{lem:anchor-annulus-Green-graphs}
There are constants \(M_\lambda<\infty\) and
\begin{equation}
 M_G=2,\qquad B_{\rm hl}=6,\qquad
 R_{\mathrm{hl},\delta}=c_{\rm hl}\delta^6,          \label{eq:anchor-half-line-radius}
\end{equation}
and \(c_{\rm hl}>0\), all independent of \(N\), with the following
properties.

\smallskip
\noindent
\emph{Positive side.}
Let \(\Phi_+^{\rm col}(p)\) be the exact inward hit of the local
\(W^u(E_N^+)\) at the collar scale
\cref{eq:anchor-local-manifold-scale}.  The raw RFDE
\cref{eq:anchor-raw-fold-time-rfde} has a unique forward solution
\(Z_+(s;p)\) from this history.  Before a connection condition is imposed,
this solution stays in the positive anchored channel and its depth-three
delay enlargement, crosses the positive inner section exactly once, and
has a common event-speed margin there.  It then continues through the
already proved central family to \(\Sigma_0\).  After three full delay
buffers, its recut compatible history and its inner event holonomy are
\(C^1\) in \(p=(\nu,\eta)\), with norm at most
\(C\delta^{-M_G}\).

Let \(\vartheta_+\) be the fixed channel phase function that equals the
local unstable coordinate on the anchor collar and the collective current
coordinate on the compact outer subannulus, joined on their common monotone
overlap.  With \(q_+(s)=\partial_sZ_+(s)\), define
\begin{equation}
 \ell_+(s)
 =\frac{D\vartheta_+(Z_{+,s})}
        {D\vartheta_+(Z_{+,s})q_+(s)},\qquad
 Q_+(s)=I-q_+(s)\ell_+(s).                            \label{eq:anchor-positive-event-projection}
\end{equation}
The denominator is nonzero by the collar and first-exit ledgers below, and
\(\ell_+(s)q_+(s)=1\) exactly.  The event-normal variational quotient, not
the full variational equation, has a causal current-path Green kernel
\(G_{+,\perp}^s(s,u):\mathbb R^{N+1}\to Q_+(s)\mathscr Y_N\).  On every fixed
positive outer subannulus,
\begin{equation}
 \|G_{+,\perp}^s(s,u)\|_{\sharp}
 \le C\delta^{-M_G}
       e^{-c_G\kappa_\delta(s-u)},\qquad s\ge u,       \label{eq:anchor-positive-causal-Green}
\end{equation}
where endpoint event correction removes the time tangent.  No decay is
claimed for the full variational process.
where, for \(\theta_m>0\),
\begin{equation}
 \kappa_\delta
 =\frac{\vartheta_0\log(1/\delta)}{\theta_m},
 \qquad
 \|\phi\|_\delta^\sharp
 =\sup_{-\theta_m\le\vartheta\le0}
       e^{2\kappa_\delta\vartheta}\|\phi(\vartheta)\|_N,
                                                               \label{eq:anchor-annulus-fading-norm}
\end{equation}
with \(0<\vartheta_0<1/8\).  If \(\theta_m=0\), take
\(\kappa_\delta=\vartheta_0a_*/\delta\).  In the delayed case,
\begin{equation}
 \|\phi\|_\delta^\sharp\le\|\phi\|_\infty
 \le\delta^{-2\vartheta_0}\|\phi\|_\delta^\sharp.     \label{eq:anchor-annulus-norm-equivalence}
\end{equation}

\smallskip
\noindent
\emph{Negative reference problem and exact chart center.}
The seed is now fixed explicitly.  Fix
\(0<\epsilon_{\rm seed}<r_{\rm seed}/2\) and work on the buffered interval
\[
 I_-^{\rm buf}=[-r_A,-r_{\rm seed}+\epsilon_{\rm seed}],
\]
which stays a fixed distance from the fold.  On this interval put
\begin{equation}
 \mathscr B_N(r)(Z,Q)
 =\mathcal J_N^{\rm fv}(r)Z-Q\mathbf1,\qquad
 (Z,Q)\in E_N\times\mathbb R .                       \label{eq:anchor-seed-normal-block}
\end{equation}
The critical-curve equations and
\(\inf_{I_-^{\rm buf}}|w_{0,N}'|>0\) give
\begin{equation}
 \sup_{N,r\in I_-^{\rm buf}}
 \|\mathscr B_N(r)^{-1}\|_{\mathbb R^N\to E_N\times\mathbb R}
 \le C_B .                                           \label{eq:anchor-seed-normal-inverse}
\end{equation}
Set \(J_*=10\).  Starting with
\(Z_0=Q_0=0\) and
\(q_0(r)=r g_{\rm anc}(r)/w_{0,N}'(r)\), define
\begin{align}
 V^{[j]}(r)&=V_{0,N}(r)+\sum_{a=1}^j\delta^aZ_a(r),\notag\\
 W^{[j]}(r)&=w_{0,N}(r)+\sum_{a=1}^j\delta^aQ_a(r),\notag\\
 q^{[j]}(r)&=q_0(r)+\sum_{a=1}^j\delta^aq_a(r).       \label{eq:anchor-seed-truncations}
\end{align}
Here \(Z_a(r)\in E_N\), \(Q_a(r),q_a(r)\in\mathbb R\), and none of these
coefficient functions depends on \(\delta\).  Write
\(j^\flat=\max\{0,j-2\}\), and let
\(\varphi_{q^{[j]}}^\tau(r)\) denote the local flow of
\(\dot r=q^{[j]}(r)\) in its time variable \(\tau\).  Define the literal
backward delay foot at fixed current by
\begin{equation}
 \Theta_k^{[j]}(r;\delta)
 =\varphi_{q^{[j]}}^{-\delta\theta_k}(r),
 \qquad 0\le k\le m .                                \label{eq:anchor-seed-fixed-current-foot}
\end{equation}
After decreasing \(\delta_0\), every such foot used below remains in
\(I_-^{\rm buf}\).  For an expression \(H(r;\delta,p)\), the notation
\begin{equation}
 [\delta^j]_rH
 :=\frac1{j!}\,\partial_\delta^jH(r;0,p)              \label{eq:anchor-seed-fixed-current-coefficient}
\end{equation}
means coefficient extraction with \(r\) and \(p\) held fixed.  For
\(1\le j\le J_*\), choose \((Z_j,Q_j)\) as the unique solution of
\begin{multline}
 [\delta^j]_r\biggl[
 \delta^2\partial_rV^{[j]}(r)q^{[j^\flat]}(r)
 -\mathscr V_N(V^{[j]}(r),W^{[j]}(r))\\
 {}-\delta^2K\sum_{k=0}^mL_{k,N}(\eta)
 \{V^{[j]}(r)-V^{[j]}(\Theta_k^{[j^\flat]}(r;\delta))\}
 \biggr]=0,                                          \label{eq:anchor-seed-voltage-recursion}
\end{multline}
and then choose \(q_j\) from
\begin{equation}
 [\delta^j]_r\left[
 \partial_rW^{[j]}(r)q^{[j]}(r)
 -(r-\delta^2\nu)g_{\rm anc}(r)
 \right]=0.                                          \label{eq:anchor-seed-resource-recursion}
\end{equation}
In \cref{eq:anchor-seed-voltage-recursion} the only order-\(j\) unknown is
acted on by \(-\mathscr B_N(r)\); after it is found,
\cref{eq:anchor-seed-resource-recursion} has the nonzero coefficient
\(w_{0,N}'(r)q_j\).  These two displayed rows are therefore a genuine
triangular recursion, not a formal appeal to an unspecified corrector.
They also give \(Z_j(-r_A)=Q_j(-r_A)=q_j(-r_A)=0\).

Let \(r_\Gamma\) solve
\begin{equation}
 r_\Gamma'=\delta q^{[J_*]}(r_\Gamma),\qquad
 r_\Gamma(0)=-r_{\rm seed},                          \label{eq:anchor-seed-slow-path}
\end{equation}
on \([ -\theta_m,\infty)\), and define
\begin{equation}
 \Gamma_-(s;p)
 =\bigl(V^{[J_*]}(r_\Gamma(s)),W^{[J_*]}(r_\Gamma(s))\bigr).
                                                               \label{eq:anchor-explicit-time-seed}
\end{equation}
The simple zero of \(q^{[J_*]}\) at \(-r_A\) implies
\(r_\Gamma(s)\to-r_A\), so this is a \(C^1\) family in
\(\mathscr D_-^0\), starts on \(\Sigma_-^{\rm seed}\), and converges to
\(E_N^-\).  It is only an approximate half-line path.  Its literal residual
\begin{equation}
 R_\Gamma(s;p)
 =\Gamma_-'(s;p)
  -\mathscr F^{\rm f}_{N,\delta,p}(\Gamma_{-,s}(p))   \label{eq:anchor-seed-residual}
\end{equation}
is computed with \cref{eq:anchor-raw-fold-time-rfde}.  Taylor's formula in
the two recursions, including the literal old-history interval
\(-\theta_m\le s-\theta_k<0\), gives
\begin{equation}
 \|R_\Gamma\|_{\mathscr F_\beta}
 +\|D_pR_\Gamma\|_{\mathscr F_\beta}
 \le C_R\delta^9
 \le\frac{c_{\rm hl}}{12C_G}
       \delta^{M_G}R_{\mathrm{hl},\delta}.             \label{eq:anchor-seed-residual-smallness}
\end{equation}
More precisely, for \(s\ge0\),
\begin{align}
 |r_\Gamma(s)+r_A|&\le Ce^{-c_q\delta s},\notag\\
 |D_pr_\Gamma(s)|&\le C(1+\delta s)e^{-c_q\delta s},               \label{eq:anchor-seed-tail-decay}\\
 \|R_\Gamma(s)\|+\|D_pR_\Gamma(s)\|
 &\le C\delta^{10}(1+\delta s)e^{-c_q\delta s}.       \label{eq:anchor-seed-tail-residual}
\end{align}
The weight is chosen with
\(\beta_\delta\le\min\{c_s,c_q\}\delta/4\), so these pointwise bounds
imply the first inequality in
\cref{eq:anchor-seed-residual-smallness}.  On the old-history interval,
\(\sup_{-\theta_m\le s\le0}|r_\Gamma(s)+r_{\rm seed}|
\le C\delta\theta_m<\epsilon_{\rm seed}\), so every prescribed old-history
value lies in \(I_-^{\rm buf}\); the coefficient recursion is never
evaluated outside its declared domain.
Here \(C_G\) is the Green constant below and the range norm is
\begin{equation}
 \|f\|_{\mathscr F_\beta}
 =\sup_{s\ge0}e^{\beta_\delta s}\delta^{-1}
       \{\|f_v(s)\|_N+|f_w(s)|\},
 \qquad 0<\beta_\delta\le\frac14c_s\delta.            \label{eq:anchor-half-line-range-norm}
\end{equation}
The order ten construction uses at most eleven ordinary derivatives of the
coefficients, including one parameter derivative.  It therefore stays
strictly inside the declared \(C^{12}\) anchor budget.  No splice to an
independently chosen local tail occurs.

The variational equation along \(\Gamma_-\) has a reference dichotomy on
\([0,\infty)\) with projections
\(P_\Gamma^s(s),P_\Gamma^u(s)\), where
\(\dim\Ran P_\Gamma^u(s)=1\).  Its forced
current-path kernels
\[
 G_{\Gamma}^s(s,u),\quad G_{\Gamma}^u(s,u):
       \mathbb R^{N+1}\longrightarrow\mathscr Y_N
\]
are obtained by variation of constants for the current path followed by
sliding-history reassembly.  They are not defined through a fictitious
bounded injection of current forcing into \(C\)-history space.  With
\begin{equation}
 \|z\|_{\mathscr Z_\beta}
 =\sup_{s\ge0}e^{\beta_\delta s}
       \|z_s\|_{s,\delta}^{\rm ph},                   \label{eq:anchor-half-line-solution-norm}
\end{equation}
where the displayed phase/fading norms are polynomially equivalent to the
ambient history norm, first put
\begin{equation}
 (\widetilde{\mathbb G}_\Gamma f)(s)
 =\int_0^sG_\Gamma^s(s,u)f(u)\,\dd u
  -\int_s^\infty G_\Gamma^u(s,u)f(u)\,\dd u.          \label{eq:anchor-reference-mixed-Green}
\end{equation}
Choose
\(e_\Gamma^{\rm ph}\in\Ran P_\Gamma^s(0)\) with
\(\ell_{\rm seed}e_\Gamma^{\rm ph}=1\), and let
\(U_\Gamma^s(s,0)e_\Gamma^{\rm ph}\) be its stable evolution.  The
phase-bordered reference Green operator is
\begin{equation}
 (\mathbb G_\Gamma^{\rm ent}f)(s)
 =(\widetilde{\mathbb G}_\Gamma f)(s)
 -U_\Gamma^s(s,0)e_\Gamma^{\rm ph}\,
      \ell_{\rm seed}(\widetilde{\mathbb G}_\Gamma f)(0).        \label{eq:anchor-reference-phase-border}
\end{equation}
It solves the same forced equation, has zero entry coordinate in
\(\Ran P_\Gamma^s(0)\cap\ker\ell_{\rm seed}\), and satisfies
\(\ell_{\rm seed}(\mathbb G_\Gamma^{\rm ent}f)(0)=0\).  Thus its initial
correction is tangent to \(\Sigma_-^{\rm seed}\); this border is imposed
before the exact base orbit is available.
Then
\begin{equation}
 \|\mathbb G_\Gamma^{\rm ent} f\|_{\mathscr Z_\beta}
 \le C_G\delta^{-M_G}\|f\|_{\mathscr F_\beta}.        \label{eq:anchor-reference-Green-bound}
\end{equation}
The second integral uses only the inverse on the transported
one-dimensional unstable bundle.

Let \(\mathcal N_\Gamma(z)\) be the exact nonlinear remainder about the
seed.  With zero stable entry coordinate in
\(\Ran P_\Gamma^s(0)\cap\ker\ell_{\rm seed}\), the contraction
\begin{equation}
 \bar z
 =\mathbb G_\Gamma^{\rm ent}
       \{-R_\Gamma+\mathcal N_\Gamma(\bar z)\}         \label{eq:anchor-chart-center-fixed-point}
\end{equation}
has a unique solution in
\(\|\bar z\|_{\mathscr Z_\beta}\le R_{\mathrm{hl},\delta}/3\).
Consequently
\begin{equation}
 \widehat Z_-(s;p)=\Gamma_-(s;p)+\bar z(s;p)
                                                               \label{eq:anchor-exact-seed-orbit}
\end{equation}
solves the raw anchored RFDE exactly, starts in
\(\Sigma_-^{\rm seed}\), and converges to \(E_N^-\).
It crosses \(\Sigma_-^{\rm ent}\) once at a \(C^1\) time
\(s_{\rm ent}(p)>3\theta_m\).  Reindex at that generated event:
\begin{equation}
 \bar Z_-(s;p)=\widehat Z_-(s+s_{\rm ent}(p);p),
 \qquad s\ge-\theta_m .                              \label{eq:anchor-exact-chart-center-orbit}
\end{equation}
This exact orbit starts in the compatible \(C^3\) part of
\(\Sigma_-^{\rm ent}\) and converges to \(E_N^-\).  It is the selected
chart-center orbit.  No invariant family is inferred from that selection.
Writing \(\widehat Z_-=(\widehat v_-,\widehat w_-)\), its entry time-tangent
history is the literal generated history
\begin{equation}
 \tau_{\rm ent}
 =\left(\vartheta\longmapsto
       \partial_s\widehat v_-(s_{\rm ent}+\vartheta),\
       \partial_s\widehat w_-(s_{\rm ent})\right),
 \qquad-\theta_m\le\vartheta\le0.                    \label{eq:anchor-compatible-entry-tangent}
\end{equation}
Because \(s_{\rm ent}>3\theta_m\), every term in this row and its first two
seam derivatives is generated by the RFDE.  Hence
\(\tau_{\rm ent}\) belongs to the tangent space of the compatible
\(C^3\) solution-manifold chart at \(\bar Z_-(0)\).  No time tangent is
used in the preceding weak seed-space phase border.
Its first hit of the local stable chart is, by definition,
\(C^-_{A,N,\delta,p}\); convergence and local stable-set identification
put that hit in \(W^s_{\rm loc}(E_N^-)\).

\smallskip
\noindent
\emph{Exact dichotomy and intrinsic stable-entry graph.}
Let \(U_-(s,u)\) be the variational solution process along
\(\bar Z_-\).  It has invariant projections \(P_-^s(s),P_-^u(s)\) on
\(\mathscr Y_N\), with \(\dim\Ran P_-^u(s)=1\), satisfying
\begin{align}
 \|U_-(s,u)P_-^s(u)\|
 &\le C\delta^{-M_G}e^{-c_s\delta(s-u)},&&s\ge u,\notag\\
 \|U_-(s,u)P_-^u(u)\|
 &\le C\delta^{-M_G}e^{-\mathcal A_u(u,s)},&&0\le s\le u,
                                                               \label{eq:anchor-annulus-Green-kernels}\\
 \mathcal A_u(u,s)&\ge c_u(u-s)/\delta
 \quad\hbox{on each fixed negative outer subannulus}. \notag
\end{align}
The second line again denotes only the inverse on the unstable line.
The projections, kernels, and their first parameter derivatives have one
common polynomial bound.

This is not inferred from the raw frozen voltage--resource system.  On
every fixed outer part ending before the anchor, parameterize the exact
orbit by its monotone collective coordinate and write
\[
 \bar Z_-(s)=\mathcal H_{\bar V,\bar w,\bar q}(\bar r(s)),
 \qquad \bar r'=\delta\bar q(\bar r).
\]
The exact resource-gauge coordinates of
\cref{prop:exact-resource-gauge-quotient} give
\begin{equation}
 (\xi_s,\omega)
 =(\psi_s,0)+a(s)E_s^c,\qquad
 a'=\delta\bar q'(\bar r)a
       +\frac{\delta}{\bar w'(\bar r)}\pi_N^Tu .       \label{eq:anchor-exact-resource-phase-splitting}
\end{equation}
The quotient history \(\psi\) has one strong-unstable line and a stable
current complement; \(aE^c\) is the longitudinal slow-stable column.
The Volterra history conjugacy in that proposition and its inverse are
uniformly bounded there because the interval stays away from the fold and
from the zero of \(\bar q\).  On the final anchor tail,
\cref{eq:anchor-full-linear-trichotomy} and asymptotic dichotomy roughness
give the same slow-stable/strong-unstable splitting; uniqueness on the
overlap identifies the bundles.  Thus the quotient dichotomy plus the slow
column is equivalent to the full \(\mathscr Y_N\) dichotomy displayed
above.  Intersecting with the
transverse entry section removes the longitudinal phase degree of freedom.
In particular, the extra frozen resource instability identified in
\cref{rem:frozen-resource-phase-quotient} is not being discarded.

Let \(\mathcal M_{{\rm ent},p}^2\) be the compatible \(C^2\)
solution-manifold chart at \(\bar Z_-(0;p)\), with the two seam rows obtained
from \cref{eq:anchor-raw-fold-time-rfde} and its first time derivative.
The base history is in its compatible \(C^3\) part by
\cref{eq:anchor-compatible-entry-tangent}.  We use the \(C^2\) graph norm
on its tangent space and the ambient fading \(C\)-norm on the range of all
half-line Green maps.  Sliding-history evolution across the three reserved
buffers maps the weak dichotomy bundles into this tangent space.  The graph
projections restrict to the generated \(C^2\) subspace, with a fixed
polynomial strong-to-weak loss; this follows directly by differentiating the
raw RFDE twice on the last two buffers, not by differentiating a translation
on arbitrary \(C\)-histories.

Here is a fixed coordinate realization of that compatible chart.  For
\(\theta_m>0\), choose
\(0<\theta_\dagger<\min\{\theta_k:\theta_k>0\}\), put
\(x=(\vartheta+\theta_\dagger)/\theta_\dagger\), and define on
\(-\theta_\dagger\le\vartheta\le0\)
\begin{align*}
 H_0(x)&=10x^3-15x^4+6x^5,\\
 H_1(x)&=-4x^3+7x^4-3x^5,\\
 H_2(x)&=\tfrac12(x^3-2x^4+x^5).
\end{align*}
Extend these functions by zero on
\([-\theta_m,-\theta_\dagger]\).  The Hermite seam lift
\begin{equation}
 \mathfrak H_\dagger(y_0,y_1,y_2)(\vartheta)
 =H_0(x)y_0+\theta_\dagger H_1(x)y_1
       +\theta_\dagger^2H_2(x)y_2                 \label{eq:anchor-entry-Hermite-seam-lift}
\end{equation}
is \(C^2\), is supported in the last buffer, and has endpoint jets
\((y_0,y_1,y_2)\).  On
\(C^2([-\theta_m,0];\mathbb R^N)\times\mathbb R\), define the section--seam
map
\begin{equation}
 \mathfrak C_{{\rm seam},p}(\phi,w)
 =\begin{pmatrix}
   \pi_N^T\{\phi(0)-\mathbf1\}+r_{\rm in}\\
   \phi'(0)-\mathscr F_v^{\rm f}(\phi,w)\\
   \phi''(0)-D\mathscr F_v^{\rm f}(\phi,w)
       [\phi',\delta\sigma_{\rm anc}(r(\phi))]
  \end{pmatrix}.                                    \label{eq:anchor-entry-section-seam-map}
\end{equation}
Its derivative on the lift
\(\mathfrak H_\dagger(a\mathbf1,y_1,y_2)\) is block lower triangular:
the three diagonal maps are \(a\mapsto a\), \(y_1\mapsto y_1\), and
\(y_2\mapsto y_2\).  Solving those rows in this order gives a fixed
right inverse \(\mathfrak R_{\rm seam}\), with a common polynomial bound,
at the base parameter.  The same triangular recursion gives
\(\mathfrak R_{{\rm seam},p}\) after shrinking the parameter chart.  Put
\begin{equation}
 \mathscr N_{\rm seam}
 =\operatorname{Ran}\mathfrak R_{\rm seam},\qquad
 \mathscr T_{{\rm ent},p}
 =\ker D\mathfrak C_{{\rm seam},p}(\bar Z_-(0;p)).
                                                               \label{eq:anchor-entry-explicit-complement}
\end{equation}
Thus the strong section tangent is explicitly the graph of
\(\mathscr T_{{\rm ent},p}\) over the fixed complement
\(\mathscr N_{\rm seam}\).  When \(\theta_m=0\), the same construction is
the ordinary finite-dimensional section chart and no seam lift is needed.

Let \(\ell_{\rm ph}\) define \(\Sigma_-^{\rm ent}\).  The exact time
tangent belongs to the stable bundle because
\(\partial_s\bar Z_-(s)\to0\).  Put
\begin{equation}
 e_-^{\rm ph}
 =\frac{\tau_{\rm ent}}
        {\ell_{\rm ph}\tau_{\rm ent}},
 \qquad
 \jmath_{\rm ent}^u
 =h_-^u-\ell_{\rm ph}(h_-^u)e_-^{\rm ph},             \label{eq:anchor-exact-phase-projected-column}
\end{equation}
Here \(h_-^u\) is normalized in the fixed entry-history norm used in
\cref{eq:anchor-half-line-radius}; this normalization keeps the
radius-\(R_{\mathrm{hl},\delta}\) nonlinear chart and its contraction
unchanged.  Thus
\(\jmath_{\rm ent}^u\) is the exact unstable
history projected into the section along an exact stable column.  Then
\begin{equation}
 \begin{aligned}
 \mathcal K_{\rm ent}(p)
 &=T_{\bar Z_-(0;p)}\mathcal M_{{\rm ent},p}^2
   \cap T_{\bar Z_-(0;p)}\Sigma_-^{\rm ent}
   \cap\Ran P_-^s(0),\\
 T_{\bar Z_-(0;p)}\mathcal M_{{\rm ent},p}^2
   \cap T_{\bar Z_-(0;p)}\Sigma_-^{\rm ent}
 &=\mathcal K_{\rm ent}\oplus
       \operatorname{span}\{\jmath_{\rm ent}^u\}.
 \end{aligned}                                         \label{eq:anchor-entry-splitting}
\end{equation}
The intersections and coordinate projections have norm at most a fixed
power of \(\delta^{-1}\).  Let
\(\Pi_{{\rm ent},p}^{s,\Sigma}\) be the projection of the strong section
tangent in \cref{eq:anchor-entry-splitting} onto \(\mathcal K_{\rm ent}(p)\)
along \(\jmath_{\rm ent}^u(p)\).  After fixing a base parameter \(p_0\), set
\begin{equation}
 \mathfrak J_p^s
 =\Pi_{{\rm ent},p}^{s,\Sigma}\big|_{\mathcal K_{\rm ent}(p_0)},
 \qquad
 \mathfrak J_p^u a=\jmath_{\rm ent}^u(p)a .           \label{eq:anchor-entry-fixed-trivialization}
\end{equation}
The first map is an isomorphism after shrinking the parameter chart; these
two explicit graph maps trivialize the domain before parameter
differentiation.  In the compatible solution-manifold entry chart,
\begin{equation}
 \Phi_{\rm ent}(\psi,u;p)
 =\bar Z_-(0;p)+\mathfrak J_p^s\psi
  +\jmath_{\rm ent}^u(p)u+\mathcal Q_{\rm ent}(\psi,u;p),
                                                               \label{eq:anchor-compatible-entry-chart}
\end{equation}
where \(\mathcal Q_{\rm ent}(\psi,u;p)\in\mathscr N_{\rm seam}\) is the
unique solution of
\[
 \mathfrak C_{{\rm seam},p}
   \bigl(\bar Z_-(0;p)+\mathfrak J_p^s\psi
       +\jmath_{\rm ent}^u(p)u+\mathcal Q_{\rm ent}(\psi,u;p)\bigr)=0.
\]
The tangent terms lie in \(\mathscr T_{{\rm ent},p}\), so
\(\mathcal Q_{\rm ent}(0,0;p)=0\),
\(D_{(\psi,u)}\mathcal Q_{\rm ent}(0,0;p)=0\), and the parameterized
implicit-function theorem gives
\(\mathcal Q_{\rm ent}=O((\|\psi\|+|u|)^2)\) with the stated fixed-chart
first-parameter bounds.  This formula specifies the complement and the
coordinate correction rather than invoking an unspecified compatible
Banach chart.

Recenter the raw RFDE at the exact orbit \(\bar Z_-\), and denote its
nonlinear remainder by \(\mathcal N_-(s,z;p)\).  Write
\begin{align}
 (\widetilde{\mathbb G}_-f)(s)
 &=\int_0^sG_-^s(s,u)f(u)\,\dd u
   -\int_s^\infty G_-^u(s,u)f(u)\,\dd u,\notag\\
 U_-^s(s,0)&=U_-(s,0)P_-^s(0),                       \label{eq:anchor-exact-mixed-Green}
\end{align}
and abbreviate the chart displacement in
\cref{eq:anchor-compatible-entry-chart} by
\(\mathfrak C_p(\psi,u)\).  For
\(\psi\in\mathcal K_{\rm ent}(p_0)\), solve the coupled infinite-horizon
system
\begin{align}
 z(s)
 &=U_-^s(s,0)P_-^s(0)\mathfrak C_p(\psi,u)
   +(\widetilde{\mathbb G}_-
          \mathcal N_-(\cdot,z;p))(s),\notag\\
 P_-^u(0)\mathfrak C_p(\psi,u)
 &=-\int_0^\infty G_-^u(0,t)
          \mathcal N_-(t,z;p)\,\dd t .               \label{eq:anchor-annulus-nonlinear-map}
\end{align}
The second row is the entry boundary row: it makes the value of the first
row at \(s=0\) equal to the compatible section history
\(\mathfrak C_p(\psi,u)\).  Since
\(P_-^u(0)D_u\mathfrak C_p(0,0)\) is an isomorphism onto the unstable
line, the two rows are a contraction for \((z,u)\).  The fold-time vector
field and the extra \(\delta^{-1}\) in
\cref{eq:anchor-half-line-range-norm} give
\begin{equation}
 \operatorname{Lip}_{\mathscr Z_\beta\to\mathscr F_\beta}
       (\mathcal N_-;R_{\mathrm{hl},\delta})
 \le C\delta^{-2}R_{\mathrm{hl},\delta},\qquad
 \|\widetilde{\mathbb G}_-\|
       \operatorname{Lip}_z(\mathcal N_-;R_{\mathrm{hl},\delta})
 \le Cc_{\rm hl}\delta^{B_{\rm hl}-M_G-2}
 =Cc_{\rm hl}\delta^2
 \le\frac13.                                         \label{eq:anchor-half-line-contraction-ledger}
\end{equation}
The second row uniquely gives the entry graph
\(u=F^{-,\rm ent}_{N,\delta,p}(\psi)\).  Hence
\begin{equation}
 \mathcal S_-^{\rm ent}(p)
 =\{\Phi_{\rm ent}(\psi,F^{-,\rm ent}_{N,\delta,p}(\psi);p):
                         \|\psi\|<R_{\mathrm{hl},\delta}\}. \label{eq:anchor-intrinsic-entry-sheet}
\end{equation}
A zero displacement is the exact chart-center orbit, so
\(F^{-,\rm ent}_{N,\delta,p}(0)=0\).
A history in this graph has a forward orbit converging to \(E_N^-\).
Conversely, every history in the declared tube whose forward orbit
converges to \(E_N^-\) eventually enters the local stable-manifold ball of
\cref{lem:anchor-scaled-local-manifolds}.  The local stable-set theorem and
\cref{eq:anchor-full-linear-trichotomy} then give decay at rate at least
\(c_s\delta/2>\beta_\delta\).  Its displacement therefore belongs to
\(\mathscr Z_\beta\).  Variation of constants on \([0,S]\), followed by
\(S\to\infty\) on the one-dimensional unstable component, gives exactly
the future integral in \cref{eq:anchor-annulus-nonlinear-map}; the terminal
term vanishes by that decay.  It hence satisfies the same weighted mixed
equation and belongs to this graph by uniqueness.  Thus the graph is
intrinsic and is independent of the approximate seed and the chosen chart
center.

The graph is \(C^2\) in \(\psi\) and \(C^1\) in \(p\).  In the fixed
parameter chart put \(x=(z,u)\),
\(\mathcal A=D_x\mathcal T_{\psi,p}(x)\), and
\(\mathcal B=D_\psi\mathcal T_{\psi,p}(x)\).  Its required rectangular
jets are the unique solutions of
\begin{align}
 (I-\mathcal A)x_p&=\partial_p\mathcal T_{\psi,p}(x),\notag\\
 (I-\mathcal A)x_\psi&=\mathcal B,\notag\\
 (I-\mathcal A)x_{\psi p}
 &=\mathcal B_p+\mathcal A_p x_\psi,
 \qquad
 \|(I-\mathcal A)^{-1}\|\le\frac32 .                \label{eq:anchor-annulus-jet-equations}
\end{align}
Here \(\mathcal A_p\) and \(\mathcal B_p\) are total first parameter
derivatives, including the already solved \(x_p\).  Thus the last row uses
only the first parameter derivative of the kernels and the second state
derivative of the nonlinear remainder; no second parameter derivative is
present.
Here \(\mathcal T_{\psi,p}\) denotes the coupled right-hand side of
\cref{eq:anchor-annulus-nonlinear-map}.  Let
\(\Pi_{\rm ent,p}=(\Pi_{\psi,p},\Pi_{u,p})\) be the local inverse-coordinate
map of \(\Phi_{\rm ent}\), and define on its physical-history image
\begin{equation}
 \mathfrak m_{\rm ent,p}(\phi)
 =\Pi_{u,p}\phi
   -F^{-,\rm ent}_{N,\delta,p}(\Pi_{\psi,p}\phi).
                                                               \label{eq:anchor-entry-membership-function}
\end{equation}
This scalar vanishes exactly on \(\mathcal S_-^{\rm ent}(p)\).  Its
normalized differential at the graph point with coordinate \(\psi\) is the
physical-history covector
\begin{align}
 \lambda_{\rm ent}(\psi;p)
 &=\{\dd u-D_\psi F^{-,\rm ent}_{N,\delta,p}(\psi)\}
     \circ D\Phi_{\rm ent}
       (\psi,F^{-,\rm ent}_{N,\delta,p}(\psi);p)^{-1},\notag\\
 \lambda_{\rm ent}(\psi;p)&
   D_u\Phi_{\rm ent}
       (\psi,F^{-,\rm ent}_{N,\delta,p}(\psi);p)1=1.   \label{eq:anchor-dual-row-pullback}
\end{align}
At \(\psi=0\), this normalization reads
\(\lambda_{\rm ent}(0;p)(\jmath_{\rm ent}^u(p))=1\).
Moreover,
\begin{equation}
 \sup_{\|\psi\|<R_{\mathrm{hl},\delta}}
 \{\|\lambda_{\rm ent}(\psi;\cdot)\|
       +\|D_p\lambda_{\rm ent}(\psi;\cdot)\|\}
 \le C\delta^{-M_\lambda}.                           \label{eq:anchor-half-line-conormal-bound}
\end{equation}
This row is obtained by differentiating the infinite mixed equation; there
is no finite terminal pullback denominator.

\smallskip
\noindent
\emph{Finite-core tail replacement and flat forgetting.}
This comparison is internal to one anchored RFDE.  Choose
\(L_\delta>0\) on the future part of the negative channel, after the common
entry and before the anchor collar, so that all delayed ancestors on
\([0,L_\delta]\) are generated by the same equation.  The comparison
section lies in a fixed compact subannulus on which
\(\bar r'=\delta\bar q(\bar r)\) and
\(0<q_-\le|\bar q|\le q_+\).  Hence
\begin{equation}
 \frac{c_L}{\delta}\le L_\delta\le\frac{C_L}{\delta},
 \qquad
 \mathcal A_u(L_\delta,0)\ge\frac{c_uL_\delta}{\delta}.
                                                        \label{eq:anchor-tail-buffer-time}
\end{equation}
Consequently
\begin{equation}
 \mathcal A_{\rm cap}(L_\delta,0)
 :=\min\{\mathcal A_u(L_\delta,0),\kappa_\delta L_\delta\}
 \ge c_{\rm cap}\frac{\log(1/\delta)}{\delta}.       \label{eq:anchor-tail-capped-action}
\end{equation}
Here \(\mathcal A_u\) is the action of the inverse on the exact
one-dimensional unstable bundle.  A finite-core tail selector is a scalar
row \(b_L(\cdot;p)\) in the compatible terminal chart at \(L_\delta\)
satisfying
\begin{equation}
 \|b_L\|_{C^2_{\rm state}C^1_p}
 +|D b_L(\jmath_L^u)|^{-1}
 \le C\delta^{-M_b},\qquad
 |D b_L(\jmath_L^u)|>0,                              \label{eq:anchor-admissible-tail-selector}
\end{equation}
whose selected terminal graph stays in the
radius-\(R_{\mathrm{hl},\delta}\) tube.
Here \(\jmath_L^u\) is the exact unstable column transported from the
entry to \(L_\delta\) and normalized in the fixed terminal chart.
This class is nonempty: the exact terminal splitting supplies the
coordinate row
\[
 b_L^{\rm can}(\phi;p)
 =\ell_L^u(p)\{\phi-\bar Z_{-,L_\delta}(p)\},
 \qquad \ell_L^u(p)\jmath_L^u(p)=1,
\]
restricted to a sufficiently small compatible terminal chart.
The entry chart, phase row, stable coordinate \(\psi\), and every
coefficient on \([0,L_\delta]\) are held fixed.  Thus a literal change of
tail coefficients or asymptotic data is supported in
\([L_\delta,\infty)\); after truncation it appears only through \(b_L\).

Eliminate \(s>L_\delta\) in the first problem by its exact
future-asymptotic tail graph and in the second by the row \(b_L=0\).
Let \(\mathscr Z_{\beta,L}\) be the restriction of
\(\mathscr Z_\beta\) to compatible current paths on \([0,L_\delta]\),
with the induced weighted norm, and put
\begin{equation}
 \mathscr W_{L,\beta}:=\mathscr Z_{\beta,L}\times\mathbb R,       \label{eq:anchor-common-core-product-space}
\end{equation}
where the scalar is the entry unstable coordinate.  Both problems are then
fixed-point maps on this same Banach space.  Let
\(Y_\infty=(z_\infty,u_\infty)\) and
\(Y_L=(z_L,u_L)\) be their fixed points, and denote the two reduced maps,
in the same trivialization, by
\(\mathcal T_\infty\) and \(\mathcal T_L\).  With
\(\Delta Y=Y_\infty-Y_L\), one has the exact value identity
\begin{align}
 \Delta Y
 &=\{I-\overline{\mathcal A}\}^{-1}
       \mathcal E_L(Y_L),\notag\\
 \overline{\mathcal A}
 &=\int_0^1D_Y\mathcal T_\infty
       (Y_L+t\Delta Y)\,\dd t,\qquad
 \mathcal E_L(Y)
 :=\mathcal T_\infty(Y)-\mathcal T_L(Y).
                                                        \label{eq:anchor-tail-value-difference}
\end{align}
For \(q=\nu\) or a direction in \(\eta\), set
\[
 \mathcal A_i=D_Y\mathcal T_i(Y_i),\qquad
 b_{i,q}=\partial_q\mathcal T_i(Y_i),\qquad
 \mathcal B_i=D_\psi\mathcal T_i(Y_i),
 \quad i\in\{\infty,L\}.
\]
Differentiating the two fixed-point equations separately and then
subtracting gives
\begin{align}
 (I-\mathcal A_\infty)\Delta Y_q
  &=(b_{\infty,q}-b_{L,q})
       +(\mathcal A_\infty-\mathcal A_L)Y_{L,q},
       &&q\in\{\nu,\eta\},                          \label{eq:anchor-tail-parameter-difference}\\
 (I-\mathcal A_\infty)\Delta(D_\psi Y)
  &=(\mathcal B_\infty-\mathcal B_L)
       +(\mathcal A_\infty-\mathcal A_L)D_\psi Y_L .
                                                        \label{eq:anchor-tail-state-difference}
\end{align}
Put \(W_i=D_\psi Y_i\), and let
\(\mathcal A_{i,p}=D_p[\mathcal A_i]\) and
\(\mathcal B_{i,p}=D_p[\mathcal B_i]\) denote total derivatives, including
the dependence through \(Y_i\).  A further differentiation, performed
separately before subtraction, gives the exact mixed equation
\begin{multline}
 (I-\mathcal A_\infty)\Delta W_p
 =(\mathcal B_{\infty,p}-\mathcal B_{L,p})
   +(\mathcal A_{\infty,p}-\mathcal A_{L,p})W_L\\
 \hspace{2.2cm}
   +\mathcal A_{\infty,p}(W_\infty-W_L)
   +(\mathcal A_\infty-\mathcal A_L)W_{L,p}.
                                                        \label{eq:anchor-tail-mixed-state-parameter-difference}
\end{multline}
These are identities in the common mixed domain.  The terms containing
\(\mathcal T_\infty-\mathcal T_L\) are precisely the future unstable
integral from \(L_\delta\), together with the fading-history tail.  If
\(\mathscr R_{\rm in}\) denotes entry evaluation followed by any fixed
inner event holonomy before the comparison buffer, then
\begin{multline}
 \bigl\|\mathscr R_{\rm in}
   (I-\overline{\mathcal A})^{-1}\mathcal E_L(Y_L)\bigr\|
 +\sum_{q\in\{\nu,\eta\}}
  \bigl\|\mathscr R_{\rm in}
    (I-\mathcal A_\infty)^{-1}
       \{(b_{\infty,q}-b_{L,q})
          +(\mathcal A_\infty-\mathcal A_L)Y_{L,q}\}\bigr\|\\
 +\bigl\|\mathscr R_{\rm in}
    (I-\mathcal A_\infty)^{-1}
       \{(\mathcal B_\infty-\mathcal B_L)
          +(\mathcal A_\infty-\mathcal A_L)D_\psi Y_L\}\bigr\|
 \le C\delta^{-M_{\rm cmp}}
       e^{-\mathcal A_{\rm cap}(L_\delta,0)}.         \label{eq:anchor-tail-Green-difference-bound}
\end{multline}
Equation
\cref{eq:anchor-tail-mixed-state-parameter-difference} gives the mixed
\(\partial_pD_\psi\) bound; it uses only the first parameter derivative of
the kernels and the \(C^2\) state derivative of the nonlinear remainder.

For clarity, the dual comparison is performed after the two graph rows
have been pulled back to the same physical entry space.  Put
\[
 F_\infty=F^{-,\rm ent}_{N,\delta,p},
 \qquad F_L=F^{-,L}_{N,\delta,p},
\]
and
\begin{align}
 \ell_i
 &=\{\dd u-D_\psi F_i(0;p)\}
       \circ D\Phi_{\rm ent}(0,F_i(0;p);p)^{-1},\notag\\
 a_i&=\ell_i(\jmath_{\rm ent}^u),\qquad
 \lambda_i=a_i^{-1}\ell_i,\qquad i\in\{\infty,L\}.
                                                        \label{eq:anchor-tail-conormal-normalization}
\end{align}
For the intrinsic row \(a_\infty=1\) at the chart center.  The finite row
is evaluated on the same fixed column; its denominator is
\(1+O(\mathfrak f_\delta)\), rather than being silently set equal to one.
The displayed normalization also covers a transported terminal row.  Its
first parameter equation is
\begin{align}
 \partial_p\lambda_i
 &=a_i^{-1}\partial_p\ell_i
   -a_i^{-2}(\partial_pa_i)\ell_i,\notag\\
 \partial_pa_i
 &=(\partial_p\ell_i)(\jmath_{\rm ent}^u)
   +\ell_i(\partial_p\jmath_{\rm ent}^u).          \label{eq:anchor-tail-conormal-parameter}
\end{align}
Consequently
\cref{eq:anchor-tail-state-difference,eq:anchor-tail-conormal-parameter}
give
\begin{equation}
 \|\lambda_\infty-\lambda_L\|
 +\|D_p(\lambda_\infty-\lambda_L)\|
 \le C\delta^{-M_{\rm cmp}}
       e^{-\mathcal A_{\rm cap}(L_\delta,0)}.         \label{eq:anchor-tail-conormal-difference}
\end{equation}

On the positive side, take two entrance representatives only after they
have entered one common physical equation.  If \(\Delta Z_+\) is their
event-aligned difference, then the exact normal equation is
\begin{equation}
 Q_+(s)\Delta Z_+(s)
 =G_{+,\perp}^s(s,s_0)Q_+(s_0)\Delta Z_+(s_0)
  +\int_{s_0}^sG_{+,\perp}^s(s,t)
       \Delta\mathcal N_+(t)\,\dd t
  +\mathcal E_{\rm ev}(s),                         \label{eq:anchor-positive-value-difference}
\end{equation}
where \(\mathcal E_{\rm ev}\) is the endpoint event correction and removes
only the time tangent.  In this display the two section times have been
frozen by the common event pullback; its moving-endpoint terms are included
in \(\mathcal E_{\rm ev}\).  Differentiating this identity once gives the
separate \(\partial_\nu\Delta Z_+\) and \(D_\eta\Delta Z_+\) equations:
\begin{align}
 \partial_q\{Q_+(s)\Delta Z_+(s)\}
={}&(\partial_qG_{+,\perp}^s)(s,s_0)
       Q_+(s_0)\Delta Z_+(s_0)\notag\\
 &+G_{+,\perp}^s(s,s_0)
       \partial_q\{Q_+(s_0)\Delta Z_+(s_0)\}\notag\\
 &+\int_{s_0}^s\bigl\{
       (\partial_qG_{+,\perp}^s)(s,t)
          \Delta\mathcal N_+(t)
       +G_{+,\perp}^s(s,t)
          \partial_q\Delta\mathcal N_+(t)\bigr\}\,\dd t\notag\\
 &+\partial_q\mathcal E_{\rm ev}(s),
 \qquad q\in\{\nu,\eta\}.                         \label{eq:anchor-positive-parameter-difference}
\end{align}
The causal estimate \cref{eq:anchor-positive-causal-Green} on a common
attracting buffer gives the same capped-action rate.

Combining the negative and positive estimates yields, for the indicated
entry and inner restrictions,
\begin{equation}
 C\delta^{-M_{\rm cmp}}
 \exp\left\{-\frac{c_{\rm anc}\log(1/\delta)}{\delta}\right\}.
                                                               \label{eq:anchor-annulus-capped-action}
\end{equation}
A change of entry stable coordinate, phase row, or compatible chart is not
a forgettable tail difference.  In particular,
\(\lambda_{\mathbf p}^{\rm ent}\) from
\cref{thm:oe-family-C1-generator-entry-conormal} is not covered merely
because it is normalized: that row is selected on a different pulled
finite core, and no nonlinear terminal graph on the exact anchored common
segment has been constructed for it.
\end{lemma}

\begin{proof}
We prove the assertions in the order in which the physical objects are
constructed.

For the positive side, use \(\Phi_+^{\rm col}\) as the prescribed old
history in \cref{eq:anchor-raw-fold-time-rfde}.  The ordinary RFDE
variation-of-constants formula separates the interval
\(0\le s<\theta_k\), where \(v(s-\theta_k)\) belongs to that prescribed
history, from \(s\ge\theta_k\), where it is generated by the current path.
The frozen current estimate
\cref{eq:frozen-current-attracting}, a continuous moving current frame,
and one Volterra--Neumann series give the event-normal quotient estimate
\cref{eq:anchor-positive-causal-Green}.  Its perturbation number is
\[
 C\delta^2\{1+\delta^{-2\vartheta_0}\}<\frac14.
\]
More explicitly, after the time column has been separated, the conjugated
normal equation has the form
\[
 y'=\mathcal A_{+,c}^0(s)y+
   \{\mathcal R_{+,{\rm fr}}+\mathcal R_{+,{\rm cor}}
      +\mathcal R_{+,{\rm del}}^{\rm gen}\}y_\bullet+f
   +\mathcal B_{+,{\rm del}}\phi_{\rm old}.
\]
The principal causal current kernel has width \(C\delta\).  In the fading
history norm the moving-frame and graph-corrector convolutions are bounded
by \(C\delta\), the generated-delay convolution by
\(C\delta^2e^{\beta_\delta\theta_m}\), and the old-history lift by
\(C\delta^2\delta^{-2\vartheta_0}\|\phi_{\rm old}\|_\infty\).  Thus
\begin{equation}
 \eta_{+,{\rm fr}}+\eta_{+,{\rm cor}}
   +\eta_{+,{\rm del}}^{\rm gen}
 \le C\{\delta+\delta^2(1+\delta^{-2\vartheta_0})\}<\tfrac14,
                                                               \label{eq:anchor-positive-Green-component-ledger}
\end{equation}
while \(\mathcal B_{+,{\rm del}}\phi_{\rm old}\) is a causal boundary
source and is not included in the operator contraction number.  Sliding the
resulting current path reassembles the literal \(C\)-history and gives
\cref{eq:anchor-positive-causal-Green}.  Differentiating this Neumann
identity once in \(p\) produces two causal current convolutions and the
derivative of the event row; the first costs \(C\delta^{-1}\) after source
normalization and the second costs the inverse section-speed factor
\(C\delta^{-1}\).  Hence the \(C^1\) kernel and event-holonomy bound is
\(C\delta^{-2}\), the same declared value \(M_G=2\) as on the negative
side.
The phase column \(q_+\) is retained exactly and is removed only by
\(Q_+\) and the endpoint event row; the proof never assigns decay to that
column.
In the scaled local unstable chart its slow coordinate \(\rho_+\) obeys
\begin{equation}
 \rho_+'=\delta\rho_+
   \{\chi_++O(\rho_+)+O(\delta)\},\qquad \chi_+\ge c_\chi>0,
 \qquad
 \|Q_+(Z_+-Z_{+,{\rm crit}})\|_{\sharp}
       \le C\delta^2|\rho_+| .                         \label{eq:anchor-positive-collar-ledger}
\end{equation}
The first row follows by dividing the exact slow eigenvalue in original
time by \(\delta\); the second is the local graph equation with the
event-normal Green cost \(O(\delta)\).  At
\(|\rho_+|=\rho_{A,\delta}\) these estimates are a fixed relative fraction
of the collar size.  They supply both the collar event-speed margin and the
initial channel bootstrap used below.
The critical current seed has a raw normal residual of size \(O(\delta)\);
the current Green convolution costs \(O(\delta)\), so the exact IVP is
\(O(\delta^2)\) from that seed in the resource-defect boundary-layer norm.
Near the local collar the exact unstable chart supplies the sharper
relative bound required to keep the solution inside the anchor interval.
Write \(q_+^{\rm red}(r)\) for the scalar reduced speed of the declared
positive reference current path, so its reference current satisfies
\(r'=\delta q_+^{\rm red}(r)\).  This scalar is distinct from the
history-valued time tangent \(q_+(s)=\partial_sZ_+(s)\) above.
Here is the first-exit ledger on the remaining compact subannulus.  Choose
the channel and its depth-three enlargement so that the reference current
path has distance at least \(d_{\rm ch}>0\) from their non-event boundary.
The causal estimate gives, up to a hypothetical first exit,
\begin{equation}
 \|Q_+(Z_+-Z_{+,\rm crit})\|_\sharp\le C\delta^2
       <\frac14d_{\rm ch},\qquad
 \max_{1\le j\le3}|r(s)-r(s-j\theta_m)|
       \le3C\delta\theta_m<\frac14d_{\rm ch}.        \label{eq:anchor-positive-first-exit-ledger}
\end{equation}
The differentiated current-row estimate in the same causal argument gives
\(|r'/\delta-q_+^{\rm red}(r)|\le C\delta\); integrating it over at most three
delay lengths gives the second estimate.  It is used only when the delayed
feet are generated; prescribed old feet were already estimated separately.
Thus neither a current state nor any of its three
delay ancestors can be the first point to leave the declared hull.  Together
with the relative collar estimate
\cref{eq:anchor-positive-collar-ledger}, this closes the first-exit
argument.  On the compact passage the reduced speed has
\(|q_+^{\rm red}|\ge q_->0\), so the same derivative estimate gives
\begin{equation}
 \frac{|r'(s)|}{\delta}\ge\frac12q_-               \label{eq:anchor-positive-event-speed-ledger}
\end{equation}
through the fixed inner section.  Monotonicity gives the unique crossing,
and the event implicit-function theorem gives its \(C^1\) dependence.  No
inverse RFDE evolution is used.

For the negative reference problem, the block
\(\mathscr B_N(r)\) is square because \(E_N\) supplies the \(N-1\) normal
voltage coordinates and \(Q\) supplies the missing collective row.
Consequently \cref{eq:anchor-seed-voltage-recursion} first determines
\((Z_j,Q_j)\); the resource row then determines \(q_j\).  Induction,
Taylor's formula with integral remainder, and
\cref{eq:anchor-seed-normal-inverse} give
\[
 \max_{j\le10}\bigl(
 \|Z_j\|_{C^{11-j}}+\|Q_j\|_{C^{11-j}}
       +\|q_j\|_{C^{11-j}}\bigr)\le C .
\]
The same induction after one \(p\)-derivative uses at most the twelfth
coefficient derivative.  The fixed-current convention in
\cref{eq:anchor-seed-fixed-current-foot,eq:anchor-seed-fixed-current-coefficient}
is essential here: it defines every Taylor coefficient before the long
half-line orbit is introduced.

To retain the weight at the asymptotic end, define the two graph residuals
at fixed \(r\) by
\begin{align}
 \mathcal E_v^{[10]}(r)
 &:=\delta^2\partial_rV^{[10]}(r)q^{[10]}(r)
   -\mathscr V_N(V^{[10]}(r),W^{[10]}(r))\notag\\
 &\quad-\delta^2K\sum_{k=0}^mL_{k,N}(\eta)
   \{V^{[10]}(r)-V^{[10]}(\Theta_k^{[10]}(r;\delta))\},\notag\\
 \mathcal E_w^{[10]}(r)
 &:=\partial_rW^{[10]}(r)q^{[10]}(r)
     -(r-\delta^2\nu)g_{\rm anc}(r).                \label{eq:anchor-seed-fixed-current-residuals}
\end{align}
Because the voltage delay row carries the prefactor \(\delta^2\), its
order-\(j\) coefficient uses only \(q_0,\ldots,q_{j-2}\).  Hence the
triangular recursion kills precisely the coefficients through order ten
of the two residuals just displayed.  Taylor's formula, also after one
\(r\)-derivative and one \(p\)-derivative, gives
\begin{equation}
 \|\partial_r\mathcal E_v^{[10]}\|
 +\|D_p\partial_r\mathcal E_v^{[10]}\|
 +|\partial_r\mathcal E_w^{[10]}|
 +|D_p\partial_r\mathcal E_w^{[10]}|
 \le C\delta^{11}.                                  \label{eq:anchor-seed-residual-r-jet}
\end{equation}
Every corrector vanishes at \(-r_A\), and every backward foot fixes
\(-r_A\).  Thus both residuals and their first \(p\)-derivatives vanish
there.  The mean-value theorem upgrades
\cref{eq:anchor-seed-residual-r-jet} to
\begin{equation}
 \|\mathcal E_v^{[10]}(r)\|+
   \|D_p\mathcal E_v^{[10]}(r)\|
 +|\mathcal E_w^{[10]}(r)|+|D_p\mathcal E_w^{[10]}(r)|
 \le C\delta^{11}|r+r_A|.                           \label{eq:anchor-seed-residual-anchor-factor}
\end{equation}

The simple stable zero gives
\(q^{[10]}(r)=-\gamma_\delta(r)(r+r_A)\), with
\(0<\gamma_-\le\gamma_\delta\le\gamma_+\) on the relevant tail.
Differentiating the scalar flow in \(p\), and using
\(D_pq^{[10]}(-r_A)=0\), yields
\begin{equation}
 |r_\Gamma(s)+r_A|\le Ce^{-c_q\delta s},
 \qquad
 |D_pr_\Gamma(s)|\le C(1+\delta s)e^{-c_q\delta s}.
                                                               \label{eq:proof-anchor-seed-tail-decay}
\end{equation}
Along this orbit the exact backward foot is
\(r_\Gamma(s-\theta_k)=
\Theta_k^{[10]}(r_\Gamma(s);\delta)\).  Consequently
\begin{equation}
 R_{\Gamma,v}(s)=\delta^{-1}
     \mathcal E_v^{[10]}(r_\Gamma(s)),
 \qquad
 R_{\Gamma,w}(s)=\delta
     \mathcal E_w^{[10]}(r_\Gamma(s)).               \label{eq:anchor-seed-residual-factorization}
\end{equation}
Equations
\cref{eq:anchor-seed-residual-anchor-factor,eq:proof-anchor-seed-tail-decay,eq:anchor-seed-residual-factorization}
give \cref{eq:anchor-seed-tail-residual}.  Since
\(\beta_\delta\le c_q\delta/4\), multiplication by the range weight and
its additional factor \(\delta^{-1}\) costs no power beyond
\(\delta^{-1}\).  The voltage row is therefore \(O(\delta^9)\) in
\(\mathscr F_\beta\), while the resource row is smaller.  This proves
\cref{eq:anchor-seed-residual-smallness} directly, with no tail splice.

We next prove the reference dichotomy rather than infer it from frozen
spectra.  In the moving current frame \(T_N(r_\Gamma(s))\), use the
critical tangent to eliminate the leading resource-to-current column.
More explicitly, with
\(E_{\Gamma,s}^c=\partial_r(V^{[10]},W^{[10]})(r_\Gamma(s))\), write
\begin{equation}
 (\xi_s,\omega)=(\psi_s,0)+a(s)E_{\Gamma,s}^c,\qquad
 a'=\delta(q^{[10]})'(r_\Gamma)a
       +\delta b_\Gamma(s)\pi_N^Tu+\mathcal E_\Gamma(s)(\psi,a),
                                                               \label{eq:anchor-reference-resource-phase-row}
\end{equation}
where \(b_\Gamma\) is uniformly bounded.  The accompanying Volterra
history coordinate is the finite-order version of
\cref{eq:resource-gauge-history-conjugacy}; its norm and inverse norm are
bounded on \(I_-^{\rm buf}\).  The full term
\(\delta b_\Gamma\pi_N^Tu\) is included in the principal operator: after a
slow convolution it is not a small perturbation.  The operator
\(\mathcal E_\Gamma\) contains only the uncancelled resource-gauge
remainder, whose coefficient norm is \(O(\delta^2)\), and the differentiated
seed residual, whose normalized norm is \(O(\delta^9)\) by
\cref{eq:anchor-seed-residual-smallness}.

In the frame of \cref{eq:generated-core-current-frame}, the current part of
the principal operator is the explicit nonautonomous ODE
\begin{equation}
 x'=\mathcal A_{\Gamma,c}^0(s)x,\qquad
 \mathcal A_{\Gamma,c}^0(s)
 =\delta^{-1}
   \begin{pmatrix}
    S_N^s(r_\Gamma(s))&0\\0&\lambda_{c,N}(r_\Gamma(s))
   \end{pmatrix}.                                    \label{eq:anchor-reference-principal-current-operator}
\end{equation}
Let \(U_{\Gamma,c}^{s}(s,u)\) denote its stable evolution and let
\(U_{\Gamma,c}^{u}(s,u)\) denote the inverse on its one-dimensional current
unstable line.  The moving-frame derivative, exact current correctors,
Volterra/delay rows, and residual rows are retained below in
\(\mathcal R_\Gamma\); none is silently absorbed into this diagonal ODE.
Put
\begin{equation}
 U_\Gamma^{\rm ph}(s,u)
 =\frac{q^{[10]}(r_\Gamma(s))}
        {q^{[10]}(r_\Gamma(u))},\qquad s\ge u.          \label{eq:anchor-reference-phase-evolution}
\end{equation}
The diagonal evolutions satisfy
\begin{align}
 \|U_{\Gamma,c}^{s}(s,u)\|
 &\le Ce^{-c(s-u)/\delta},&&s\ge u,\notag\\
 \|U_{\Gamma,c}^{u}(s,u)\|
 &\le Ce^{-c(u-s)/\delta},&&0\le s\le u,\notag\\
 |U_\Gamma^{\rm ph}(s,u)|
 \le Ce^{-c_s\delta(s-u)},&&s\ge u.                  \label{eq:anchor-reference-principal-evolution}
\end{align}
For the last bound, split \(I_-\) into a fixed compact part and a
neighborhood of the simple stable zero at \(-r_A\).  On the compact part
the transit time is \(O(\delta^{-1})\) and the ratio is uniformly bounded;
near the anchor the simple-zero estimate gives the displayed exponential.
This proves the slow stable block on the whole half-line.

The full principal stable evolution is lower triangular.  Its phase
component, for stable current data \(x\) and phase data \(a\), is
\begin{multline}
 U_{\Gamma,0}^{s,{\rm ph}}(s,u)(x,a)
 =U_\Gamma^{\rm ph}(s,u)a\\
 \quad+\int_u^sU_\Gamma^{\rm ph}(s,\zeta)
       \delta b_\Gamma(\zeta)\pi_N^T
       U_{\Gamma,c}^{s}(\zeta,u)x\,\dd\zeta .          \label{eq:anchor-reference-stable-offdiagonal}
\end{multline}
The fast width is \(O(\delta)\), and therefore
\begin{equation}
 \left\|\int_u^sU_\Gamma^{\rm ph}(s,\zeta)
       \delta b_\Gamma(\zeta)\pi_N^T
       U_{\Gamma,c}^{s}(\zeta,u)\,\dd\zeta\right\|
 \le C\delta^2e^{-c_s'\delta(s-u)}.                  \label{eq:anchor-reference-stable-offdiagonal-bound}
\end{equation}

The principal unstable bundle is a graph over the current unstable line.
Pull its row on the asymptotic anchor tail backward and denote the result by
\(H_\Gamma^u(0)\).  Invariance gives, for \(s\ge0\),
\begin{multline}
 H_\Gamma^u(s)=U_\Gamma^{\rm ph}(s,0)H_\Gamma^u(0)
       U_{\Gamma,c}^{u}(0,s)\\
 \quad+\int_0^sU_\Gamma^{\rm ph}(s,\zeta)
       \delta b_\Gamma(\zeta)\pi_N^T
       U_{\Gamma,c}^{u}(\zeta,s)\,\dd\zeta .          \label{eq:anchor-reference-unstable-graph-row}
\end{multline}
Both the tail row and the integral are bounded by \(C\delta^2\), since the
current inverse has width \(O(\delta)\).  Thus
\begin{equation}
 \sup_{s\ge0}\|H_\Gamma^u(s)\|\le C\delta^2,          \label{eq:anchor-reference-unstable-graph-bound}
\end{equation}
and
\begin{equation}
 U_{\Gamma,0}^{u}(s,u)
  \binom{x}{H_\Gamma^u(u)x}
 =\binom{U_{\Gamma,c}^{u}(s,u)x}
          {H_\Gamma^u(s)U_{\Gamma,c}^{u}(s,u)x},
 \qquad 0\le s\le u.                                 \label{eq:anchor-reference-full-unstable-inverse}
\end{equation}
The inverse on this graph has the same strong exponential bound.  Hence the
normal quotient has one current unstable line and the longitudinal phase is
stable; the second raw frozen resource instability in
\cref{rem:frozen-resource-phase-quotient} has not been silently discarded.

Let \(\Pi_{\Gamma,0}^{s/u}(s)\) be the resulting graph projections.  For a
conjugated current/phase forcing define
\begin{equation}
 (\mathbb G_{\Gamma,0}f)(s)
 =\int_0^sU_{\Gamma,0}^s(s,u)\Pi_{\Gamma,0}^s(u)f(u)\,\dd u
 -\int_s^\infty U_{\Gamma,0}^u(s,u)
       \Pi_{\Gamma,0}^u(u)f(u)\,\dd u .                \label{eq:anchor-reference-principal-Green}
\end{equation}
In the normalized source norm of
\cref{eq:anchor-half-line-range-norm}, the forcing has current size
\(\delta\|f\|_{\mathscr F_\beta}\).  The two current kernels have width
\(C\delta\), the direct phase kernel has width \(C\delta^{-1}\), the
stable-current-to-phase kernel has amplitude \(C\delta^2\) and slow width
\(C\delta^{-1}\), and the unstable graph has norm \(C\delta^2\).  Hence
\begin{equation}
 \|\mathbb G_{\Gamma,0}\|_{\mathscr F_\beta\to\mathscr Z_\beta}
 \le C.                                               \label{eq:anchor-reference-principal-Green-bound}
\end{equation}

Write the remaining conjugated operator as
\[
 \mathcal R_\Gamma
 =\mathcal R_{\rm fr}+\mathcal R_{\rm cor}
  +\mathcal R_{\rm rg}+\mathcal R_{\rm del}
  +\mathcal R_{\rm res}.
\]
The five terms are respectively the moving-frame row, the finite current
correctors, the resource-gauge remainder after the complete triangular
term in \cref{eq:anchor-reference-resource-phase-row} has been removed,
the literal delay translation, and the differentiated seed residual.
Convolution with \(\mathbb G_{\Gamma,0}\) gives the whole-half-line
perturbation number
\begin{equation}
 \begin{split}
 \eta_j&:=\|\mathbb G_{\Gamma,0}\mathcal R_j\|_
       {\mathscr Z_\beta\to\mathscr Z_\beta},
       \qquad j\in\{{\rm fr,cor,rg,del,res}\},\\
 \eta_\Gamma
 &\le\eta_{\rm fr}+\eta_{\rm cor}+\eta_{\rm rg}
       +\eta_{\rm del}+\eta_{\rm res},\\
 \eta_{\rm fr}&\le C\delta,\qquad
 \eta_{\rm cor}\le C\delta,\qquad
 \eta_{\rm rg}\le C\delta,\qquad
 \eta_{\rm del}\le
 C\delta^2(1+\delta^{-2\vartheta_0}),\qquad
 \eta_{\rm res}\le C\delta^7,\\
 \eta_\Gamma
 &\le C\{\delta+\delta^2(1+\delta^{-2\vartheta_0})\}
 \le\tfrac14.                                        \label{eq:anchor-reference-dichotomy-number}
 \end{split}
\end{equation}
The fast current Green width is \(O(\delta)\).  The only resource-gauge
coefficient convolved with the slow width \(O(\delta^{-1})\) is now the
uncancelled \(O(\delta^2)\) remainder, so its contribution is
\(O(\delta)\).  The full \(\delta b_\Gamma\pi_N^Tu\) term has already been
included in \(\mathbb G_{\Gamma,0}\).  This proves the first line
component by component.
Here a generated delayed value costs
\(e^{\beta_\delta\theta_m}\), whereas on
\(0\le s<\theta_k\) the prescribed seed history is estimated directly by
\(\delta^{-2\vartheta_0}\).  Thus old and generated histories are not
identified: on the first delay interval write the delay row as
\(\mathcal R_{\rm del}^{\rm gen}z+\mathcal B_{\rm del}\phi_0\), where the
first term contains only \(s-\theta_k\ge0\) and the second is a boundary
lifting of the prescribed history.  Only
\(\mathcal R_{\rm del}^{\rm gen}\) enters \(\eta_{\rm del}\); the lifting
obeys the same displayed polynomial old-history bound.  The whole-half-line
Neumann series is
\begin{equation}
 \mathbb G_\Gamma
 =(I-\mathbb G_{\Gamma,0}\mathcal R_\Gamma)^{-1}
       \mathbb G_{\Gamma,0},\qquad
 \|(I-\mathbb G_{\Gamma,0}\mathcal R_\Gamma)^{-1}\|
 \le\frac43 .                                        \label{eq:anchor-reference-Green-Neumann}
\end{equation}
It gives \(P_\Gamma^{s/u}\) and the forced current-path kernels, and
sliding-history reassembly gives
\cref{eq:anchor-reference-Green-bound}.

For a parameter direction \(\dot p\), differentiation of
\(\mathbb G_\Gamma=\mathbb G_{\Gamma,0}
 +\mathbb G_{\Gamma,0}\mathcal R_\Gamma\mathbb G_\Gamma\) gives
\begin{multline}
 D_p\mathbb G_\Gamma[\dot p]
 =(I-\mathbb G_{\Gamma,0}\mathcal R_\Gamma)^{-1}
 \bigl\{D_p\mathbb G_{\Gamma,0}[\dot p]
 +D_p\mathbb G_{\Gamma,0}[\dot p]
       \mathcal R_\Gamma\mathbb G_\Gamma\\
 \hspace{42mm}
 +\mathbb G_{\Gamma,0}D_p\mathcal R_\Gamma[\dot p]
       \mathbb G_\Gamma\bigr\}.                    \label{eq:anchor-reference-Green-first-jet}
\end{multline}
Differentiating a current kernel produces at most two fast convolutions;
differentiating \cref{eq:anchor-reference-phase-evolution} produces at
most two slow convolutions.  After the source factor \(\delta\), their
combined cost is at most \(C\delta^{-1}\).  Differentiating a graph
projection or the seed-section phase normalization may cost one further
factor \(\delta^{-1}\), because the physical section speed is of order
\(\delta\).  The first parameter derivatives of the five remainder rows
therefore satisfy
\begin{equation}
 \|D_p\mathbb G_{\Gamma,0}\|
 +\|D_p\mathbb G_{\Gamma,0}\mathcal R_\Gamma\|
 \le C\delta^{-2}.                                   \label{eq:anchor-reference-principal-first-jet-bound}
\end{equation}
The five differentiated remainder rows are not hidden in the last display:
\begin{align}
 \|\mathbb G_{\Gamma,0}D_p\mathcal R_{\rm fr}\|
 +\|\mathbb G_{\Gamma,0}D_p\mathcal R_{\rm cor}\|
 +\|\mathbb G_{\Gamma,0}D_p\mathcal R_{\rm rg}\|
 &\le C\delta^{-2},\notag\\
 \|\mathbb G_{\Gamma,0}D_p\mathcal R_{\rm del}\|
 &\le C\delta^{-2}
      \{1+\delta^{2-2\vartheta_0}\},\notag\\
 \|\mathbb G_{\Gamma,0}D_p\mathcal R_{\rm res}\|
 &\le C\delta^5.                                      \label{eq:anchor-reference-five-remainder-jets}
\end{align}
For the delay row the generated and prescribed-old-history pieces are
differentiated separately as above; the displayed bound is deliberately
weaker than their component bounds but uniform in the common parameter
norm.  Summing these five rows recovers
\(\|\mathbb G_{\Gamma,0}D_p\mathcal R_\Gamma\|\le C\delta^{-2}\).
Together with
\cref{eq:anchor-reference-Green-Neumann,eq:anchor-reference-Green-first-jet},
this proves the declared common bound \(C_G\delta^{-2}\), namely
\(M_G=2\), for the Green operator, the projections, and their first
parameter derivatives.
The same construction supplies a stable phase column
\(e_\Gamma^{\rm ph}\) transverse to the seed section, proving that
\cref{eq:anchor-reference-phase-border} is bounded.  Current forcing was
handled throughout by the displayed current-path kernels, never by a
history-space injection.

The residual bound and
\cref{eq:anchor-reference-Green-bound} put the affine term in
\cref{eq:anchor-chart-center-fixed-point} inside the
\(R_{\mathrm{hl},\delta}/12\)
ball.  Moreover,
\[
 \|\mathbb G_\Gamma^{\rm ent}\|
 \operatorname{Lip}_{\mathscr Z_\beta\to\mathscr F_\beta}
       (\mathcal N_\Gamma;R_{\mathrm{hl},\delta})
 \le Cc_{\rm hl}\delta^{6-2-2}
 =Cc_{\rm hl}\delta^2<\frac13.
\]
Thus the fixed point exists before an exact center orbit or its dichotomy
is used.  It makes \(\widehat Z_-\) an exact orbit.  The phase border makes
\(\ell_{\rm seed}\bar z(0)=0\), so this exact orbit starts on
\(\Sigma_-^{\rm seed}\).

On the fixed segment
\(-r_{\rm in}\le r\le-r_{\rm seed}\), the reference slow speed has
\(-q^{[10]}\ge c_q>0\).  The correction to the physical event speed is at
most \(C\delta^{-1}R_{\mathrm{hl},\delta}=C\delta^5\), whereas the reference speed has
size \(c_q\delta\).  Hence \(\widehat Z_-\) crosses
\(\Sigma_-^{\rm ent}\) exactly once, and
\[
 s_{\rm ent}\ge\frac{r_{\rm in}-r_{\rm seed}}{C\delta}
       >3\theta_m .
\]
The event implicit-function theorem gives its \(C^1\) parameter
dependence.  Recutting at this generated event gives
\(\bar Z_-\) in \cref{eq:anchor-exact-chart-center-orbit}.

The variational coefficients along \(\widehat Z_-\), and hence along the
reindexed \(\bar Z_-\), differ from the reference coefficients by
\(O(R_{\mathrm{hl},\delta}/\delta)\).  In the normalized source norm the corresponding
Green perturbation number is at most
\(C\delta^{-M_G-2}R_{\mathrm{hl},\delta}=C\delta^2\).  Dichotomy roughness therefore
transfers the reference splitting and Green bounds to that exact orbit.
This order removes both the base-orbit and the entry-phase circularities.

For completeness, this roughness statement is interpreted in the exact
resource-phase quotient, not in the raw frozen full system.  The
monotonicity estimate used for the entry crossing persists up to a fixed
local anchor section.  Hence \(\bar Z_-\) has the exact graph
parameterization required by
\cref{prop:exact-resource-gauge-quotient} on that outer portion.  In its
coordinates the normal RFDE is the whole-half-line Neumann perturbation
already estimated above, while the longitudinal column is the exact
tangent with cocycle
\(\bar q(\bar r(s))/\bar q(\bar r(u))\).  The bounded Volterra conjugacy
and the transverse phase row lift these two pieces to the full section.
On the final tail, asymptotic roughness of
\cref{eq:anchor-full-linear-trichotomy} gives the same bundles; overlap
uniqueness concatenates them without a terminal boundary row.  This
supplies the quotient/full equivalence that a frozen current root count
alone would not provide.

The exact time tangent is a stable solution and is transverse to
\(\Sigma_-^{\rm ent}\).  Hence
\cref{eq:anchor-exact-phase-projected-column,eq:anchor-entry-splitting}
follow from the exact one-dimensional unstable bundle.  Standard graph
formulas give the fixed
trivializations \(\mathfrak J_p^{s/u}\) with the declared polynomial
bounds.  The exact mixed equation
\cref{eq:anchor-annulus-nonlinear-map} is then a contraction by
\cref{eq:anchor-half-line-contraction-ledger}.  Its second row equates the
chart's unstable component to the infinite future integral; it is not a
hidden finite terminal row.  This proves both directions of the membership
statement without a backward RFDE semiflow.

The reference path remains a fixed distance inside the negative channel
until it enters the local anchor ball.  The base correction is at most
\(R_{\mathrm{hl},\delta}/3\), and the graph correction for
\(\|\psi\|\le R_{\mathrm{hl},\delta}\) is at most
\(CR_{\mathrm{hl},\delta}\).  The phase/fading norm
controls every literal delayed value on the depth-three hull.  Decreasing
\(c_{\rm hl}\) therefore gives negative channel containment for the base
and the whole graph before either is continued into the central chart.
Convergence places the local hit on \(W^s_{\rm loc}(E_N^-)\), which
justifies the later chart-center notation.

After three complete delay buffers, recut an exact weak path by
\begin{equation}
 \mathfrak S_3Z
 =\bigl(\vartheta\mapsto v(3\theta_m+\vartheta),
                   w(3\theta_m)\bigr).               \label{eq:anchor-three-buffer-recut}
\end{equation}
The image lies in the compatible \(C^3\) solution-manifold chart.  With
\(\mathscr F_v^{\rm f}\) denoting the voltage component of the raw field,
its three seam rows are
\begin{equation}
 \phi(0)=v(3\theta_m),\qquad
 \phi'(0)=\mathscr F_v^{\rm f}(\phi,w),\qquad
 \phi''(0)=D\mathscr F_v^{\rm f}(\phi,w)
       [\phi',\delta\sigma_{\rm anc}(r(\phi))].        \label{eq:anchor-three-seam-rows}
\end{equation}
They follow directly by differentiating
\cref{eq:anchor-raw-fold-time-rfde}; all delayed ancestors are generated
after the third buffer.  This is the
fixed strong chart at the preliminary seed flow.  The remaining regular
forward segment to \(s_{\rm ent}>3\theta_m\), followed by the fixed event
pullback, transports it to \(\mathcal M_{{\rm ent},p}^2\) and to the
explicit trivialization
\cref{eq:anchor-entry-fixed-trivialization}.  This is the chart in which the first parameter equation
\cref{eq:anchor-annulus-jet-equations} and the conormal derivative are
taken.  Differentiating the contraction once gives the \(3/2\) inverse
bound and \cref{eq:anchor-half-line-conormal-bound}.

It remains to prove the tail-replacement assertions.  Choose the comparison
section in a fixed compact subannulus of the negative channel.  Since
\(\bar r'=\delta\bar q(\bar r)\) there and \(\bar q\) is bounded above and
below away from zero, its fold-time distance from the entry is at least
\(c/\delta\).  The inverse unstable estimate and sliding-history
reassembly therefore give
\begin{equation}
 \|U_-^u(0,L_\delta)\|
 +\sup_{t\ge L_\delta}
    e^{\beta_\delta(t-L_\delta)}
      \|G_-^u(0,t)\|
 \le C\delta^{-M_G}
       e^{-\mathcal A_{\rm cap}(L_\delta,0)}.        \label{eq:anchor-tail-entry-kernel-bound}
\end{equation}
The event pullback freezes this section before parameter differentiation.
Differentiating the kernel once changes only the polynomial prefactor.

To see exactly where a terminal selector enters, solve
\(b_L=0\) for its terminal unstable coordinate
\(\gamma_L\) by
\cref{eq:anchor-admissible-tail-selector}.  It and its two state
derivatives and first parameter derivative are bounded by a fixed power of
\(\delta^{-1}\).  The entry-unstable component of the difference between
the infinite and finite mixed maps is
\begin{equation}
 \mathcal E_L^{\rm ent}(Y)
 =-\int_{L_\delta}^{\infty}
       G_-^u(0,t)\mathcal N_-
          (t,\operatorname{Ext}_\infty Y;p)\,\dd t
   -U_-^u(0,L_\delta)\jmath_L^u
       \gamma_L(Y_{L_\delta};p)
   +\mathcal E_L^{\rm coef}(Y).                    \label{eq:anchor-tail-explicit-defect}
\end{equation}
The last term is zero for a literal truncation.  More generally it is the
same future unstable integral with the coefficient difference inserted,
and its support is \(t\ge L_\delta\).  The entry-stable component is zero
because \(\psi\), the phase row, and the entry chart are fixed.  Equations
\cref{eq:anchor-tail-entry-kernel-bound,eq:anchor-tail-explicit-defect}
give the value defect in
\cref{eq:anchor-tail-Green-difference-bound}.

Subtracting
\(Y_i=\mathcal T_i(Y_i)\) and applying the mean-value formula proves
\cref{eq:anchor-tail-value-difference}.  Differentiation of each fixed-point
equation gives
\[
 (I-\mathcal A_i)Y_{i,q}=b_{i,q},\qquad
 (I-\mathcal A_i)D_\psi Y_i=\mathcal B_i .
\]
Subtracting these two pairs gives
\cref{eq:anchor-tail-parameter-difference,eq:anchor-tail-state-difference}
with their displayed types.  The contraction estimate
\cref{eq:anchor-half-line-contraction-ledger} bounds every inverse by
\(3/2\).  Differentiating
\cref{eq:anchor-tail-explicit-defect} once in \(p\), and once in \(p\)
after one state derivative, preserves its future support and costs only a
fixed power of \(\delta^{-1}\).  In particular, each difference
\(\mathcal B_{\infty,p}-\mathcal B_{L,p}\) and
\(\mathcal A_{\infty,p}-\mathcal A_{L,p}\) in
\cref{eq:anchor-tail-mixed-state-parameter-difference} is the corresponding
state derivative of the future-supported defect, plus terms already
containing \(\Delta Y\), \(\Delta Y_p\), or \(\Delta W\).  Thus
\begin{multline}
 \bigl\|\mathscr R_{\rm in}(I-\mathcal A_\infty)^{-1}
 \{(\mathcal B_{\infty,p}-\mathcal B_{L,p})
   +(\mathcal A_{\infty,p}-\mathcal A_{L,p})W_L\\
 \hspace{2.1cm}
   +\mathcal A_{\infty,p}(W_\infty-W_L)
   +(\mathcal A_\infty-\mathcal A_L)W_{L,p}\}\bigr\|
 \le C\delta^{-M_{\rm cmp}}
       e^{-\mathcal A_{\rm cap}(L_\delta,0)}.       \label{eq:anchor-tail-mixed-Green-bound}
\end{multline}
Together with
\cref{eq:anchor-tail-Green-difference-bound}, this proves the value,
the two first-parameter, the state-derivative, and the mixed
state--parameter estimates.

Here \(\operatorname{Ext}_\infty Y\) is the exact tail-graph extension
used when the infinite problem is reduced to the common core.  The chart
pullback in
\cref{eq:anchor-tail-conormal-normalization} is common to the two rows.
The normalization formula and its derivative are the ordinary quotient
rule, which gives \cref{eq:anchor-tail-conormal-parameter}.  Insert the
\(D_\psi F\) and \(\partial_pD_\psi F\) estimates just proved.  The
intrinsic denominator is one in the compatible chart.  The finite
denominator differs from it by the just-proved graph-jet bound and is
therefore bounded away from zero; the same remains true after any admitted
common transport.  This proves
\cref{eq:anchor-tail-conormal-difference}.

Finally, \cref{eq:anchor-positive-value-difference} is the ordinary
current-path variation-of-constants identity projected by \(Q_+\).
The event correction is obtained by differentiating the common section
equation, so it removes the time column and no stable current column.
The nonlinear difference in its integral is absorbed with the same
Volterra--Neumann margin used for
\cref{eq:anchor-positive-causal-Green}.  Its \(\nu\)- and \(\eta\)-derivatives
satisfy the differentiated causal identity and incur polynomial losses
only.  Since the positive comparison buffer also has fold-time length at
least \(c/\delta\), the fading cap proves
\cref{eq:anchor-annulus-capped-action}.  This completes the proof.
\end{proof}

\begin{proof}[Proof of \cref{prop:anchor-central-finite-shooting}]
The two-scale identities in
\cref{lem:anchor-reserved-physical-overlap} put the selected central trace
and the independently exact outer reference on the same generated physical
section \(\Sigma_b\).  In particular, their complete delay hulls belong to
the unmodified anchored RFDE, and
\cref{eq:anchor-reserved-component-overlap} puts them in one compatible
chart.  The
dynamic resource term \(\delta^2/(2\alpha)\) is included on both sides; thus
this comparison is not an identification with the frozen critical curve.

The same-core terminal construction
\cref{lem:anchor-same-core-intrinsic-terminal} supplies the exact
future-decaying sheet at \(L\) and the positive cross pairing on the outer
unstable column.  Using that row as the terminal membership condition, the
literal-history mixed problem in
\cref{lem:anchor-intrinsic-sheet-at-b} constructs the graph
\(\mathcal S_b^-(p)\).  Its entry row is equality in the full compatible
history space.  Hence a zero belongs to the stable set of the fixed anchored
RFDE by forward uniqueness; no finite selector, backward RFDE semiflow, or
preparation-dependent endpoint condition occurs.

At the last exact-agreement cut, the selected terminal row defines the
finite graph in \cref{lem:anchor-selected-finite-sheet-comparison}.  The
selected central trace belongs to that graph exactly, while the intrinsic
and finite graphs differ in \(C^1\), including their selected first
parameter jets, by the preassigned order
\(\varepsilon_{bc,\delta}\).  This proves
\cref{eq:anchor-main-finite-intrinsic-comparison} without asserting that
the physical stable sheet has small curvature.

The central projective coefficient is evaluated only after this intrinsic
row has been constructed.
\Cref{lem:anchor-central-intrinsic-projective-coefficient} proves
\cref{eq:anchor-main-projective-action}.  In particular, the stable-sheet
slope is compared with the selected finite row before it is applied to the
central scalar column.  The action-normalized rectangle estimate
\cref{eq:anchor-intrinsic-rectangle-ledger} therefore gives
\cref{eq:anchor-main-action-strip-containment}.  This is the radius estimate
needed for every delay hull in the shooting strip to remain inside the exact
central graph region.

Apply \cref{thm:anchor-intrinsic-Lin-intersection}.  Its oriented derivative
and opposite face signs prove existence and uniqueness before the
parameterized implicit-function theorem is used.  The same theorem gives
\cref{eq:anchor-main-intrinsic-root-distance,%
eq:anchor-main-intrinsic-root-flatness}; its literal-history zero and forward
conormal identity give all remaining stable-sheet assertions in the
proposition.

It remains only to locate the actual incoming normal coordinate.  The
positive channel has the causal event-normal Green kernel in
\cref{eq:anchor-positive-causal-Green}.  Compare the exact incoming anchor
branch with the attracting selected reference, differentiate the causal
identity \cref{eq:anchor-positive-value-difference}, and use the fixed
attracting buffer.  This gives
\begin{equation}
 \|\psi_A-\bar\psi_p\|
 \le C\delta^{-M_+}e^{-a_+/\delta}.                  \label{eq:anchor-positive-reference-flatness}
\end{equation}
The stable shooting radius in
\cref{eq:anchor-intrinsic-strip-radii} is bounded below by a fixed
algebraic power, whereas the right side above is smaller than every
algebraic power.  After decreasing \(\delta_0\), this proves
\cref{eq:anchor-incoming-in-shooting-ball}.
\end{proof}

\begin{proof}[Proof of \cref{thm:anchor-annulus-flat-forgetting}]
The positive half-line lemma gives the exact incoming hit of the inner
section using only forward RFDE evolution.  Continue it through the
attracting original-endpoint chart after it enters
\(|r|<r_{\rm loc}\).  Forward uniqueness identifies the chart segment with
the same anchored orbit and gives its exact hit \(A^0\) of \(\Sigma_0\).

On the negative side, the generated cut \(\Sigma_b\) and the forward hit
\(\mathcal P_{0\to b,p}\) lie wholly inside the exact central tube.
\Cref{lem:anchor-intrinsic-sheet-at-b} identifies
\(\mathcal S_b^-(p)\) exactly as the compatible histories whose forward
orbits remain in the declared channel and converge to \(E_N^-\).  This
description agrees, after forward evolution to the later entry section,
with the intrinsic half-line sheet in
\cref{eq:anchor-intrinsic-entry-sheet}; the equality follows from forward
uniqueness in the common anchored equation.

The raw compatible central coordinates make
\[
 \phi\longmapsto
 \bigl(\mathfrak n_p(\phi),
       c_{\mathscr H}\mathscr H_\alpha(z(\phi))\bigr)
\]
a \(C^2\)-in-state, selected-\(C^1\)-in-parameter map into
\(\mathscr K_0\times\mathbb R\).  Its differential is block triangular:
the first diagonal block is the normal/history coordinate identity, and
the scalar block is nonzero by
\cref{eq:anchor-first-integral-local-invertibility}.  This gives the exact
chart
\cref{eq:anchor-central-history-chart,eq:anchor-central-chart-identities}.

Compose the intrinsic membership at \(b\) with the forward hit and the
central chart.  The result is \(\mathcal M_{0,p}^{\rm int}\) in
\cref{eq:anchor-main-intrinsic-membership}.
\Cref{prop:anchor-central-finite-shooting} proves an opposite-sign bracket,
then a unique zero graph on a ball containing \(\psi_A\); it does not infer
nonemptiness from a conditional derivative.

At a point of this graph put
\begin{align}
 k_0^u&=D_u\mathfrak C_{0,p}(\psi,u)1,\notag\\
 \phi_b&=\mathcal P_{0\to b,p}
              (\mathfrak C_{0,p}(\psi,u)).
                                                               \label{eq:anchor-central-entry-column-data}
\end{align}
Applying the intrinsic row at \(b\) gives the exact typed identity
\begin{equation}
 \partial_u\mathcal M_{0,p}^{\rm int}(\psi,u)
 =D\mathfrak m_{b,p}(\phi_b)
      D\mathcal P_{0\to b,p}k_0^u.                  \label{eq:anchor-forward-column-coefficient-identity}
\end{equation}
The projective coefficient
\cref{lem:anchor-central-intrinsic-projective-coefficient} and the
action-balanced oscillation estimate in
\cref{thm:anchor-intrinsic-Lin-intersection} give its positive orientation
throughout the strip.  This proves
\cref{eq:anchor-forward-membership-transversality}.  The argument uses only
the annihilator of the intrinsic stable sheet; it never identifies a domain
column with a cokernel or chooses an inverse history evolution.

The graph is \(C^2\) in state and \(C^1\) in parameters by the final,
now noncircular, parameterized implicit-function step in the proposition.
This proves
\cref{eq:anchor-stable-sheet-graph,eq:anchor-central-graph-preimage-identity}.
Differentiating the composition gives the forward covector pullback
\cref{eq:anchor-central-conormal-forward-pullback}; no inverse holonomy is
used.  For an admissible finite-core tail selector, apply
\cref{eq:anchor-tail-value-difference,eq:anchor-tail-parameter-difference,%
eq:anchor-tail-state-difference} and then the common forward event
holonomies.  This gives every history and graph operator in
\cref{eq:anchor-flat-jet-list}.  The normalized row calculation
\cref{eq:anchor-tail-conormal-normalization,%
eq:anchor-tail-conormal-parameter} gives
\cref{eq:anchor-flat-conormal-comparison}.  Taking the smallest capped
action among the finitely many forward concatenations gives
\cref{eq:anchor-flat-envelope} and proves the theorem.  At no point is the
finite physical-speed row of
\cref{thm:oe-family-C1-generator-entry-conormal} used as a defining row for
the intrinsic stable sheet.
\end{proof}

\subsection{The zero fiber and its response}

\begin{lemma}[Actual central traces and normal defect]
\label[lemma]{lem:anchor-actual-central-defect}
Let \(A^0\) be the incoming hit and \(B^0\) its paired stable-sheet point
from \cref{eq:anchor-paired-stable-point}.  Their exact RFDE orbit
segments admit the raw central coordinates
\((z^\sigma,h^\sigma)\), \(\sigma\in\{a,r\}\).  On the retained delay
hull define the current normal defect and its sliding history by
\[
 \rho^\sigma(s)
 =h^\sigma(s)-H_{N,\delta,\nu,\eta}(z^\sigma(s)),
 \qquad
 \rho_s^\sigma(\vartheta)=\rho^\sigma(s+\vartheta).
\]
Put
\begin{equation}
 J_{a,S}=[-S_\delta,0],\qquad J_{r,S}=[0,S_\delta],
 \qquad
 \|\rho\|_{\mathcal X_{\sigma,S}^\sharp}
 =\sup_{s\in J_{\sigma,S}}\|\rho_s\|_\delta^\sharp .
                                                               \label{eq:anchor-retained-defect-space}
\end{equation}
These intervals, including every delay foot used in the displayed sliding
histories, lie in the uncut
retained physical tube on which the graph
\(H_{N,\delta,\nu,\eta}\) is an exact invariant graph of the unmodified
RFDE.  The preceding defect identity is used only on these intervals.  For
\(\mathfrak A_1=\{1,\partial_\nu,D_\eta\}\),
\begin{equation}
 \|\mathscr D\rho^\sigma\|_{\mathcal X_{\sigma,S}^\sharp}
 \le C\mathfrak f_\delta,\qquad
 \mathscr D\in\mathfrak A_1.                       \label{eq:anchor-actual-graph-defect-flat}
\end{equation}
The exact planar field has the solution-wise expansion
\begin{equation}
 \frac{\dd}{\dd s}z^{\sigma,\rm anc}
 =q_0+\delta q_{1,N}+\delta^2q_{2,N}
   +\delta^3\mathcal R_{3,N}
   +\mathcal B_N^\sigma\rho^\sigma
   +\mathcal N_N^\sigma(\rho^\sigma),              \label{eq:anchor-actual-field-expansion}
\end{equation}
where
\(\|\mathcal B_N^\sigma\|\le C\langle s\rangle^Me^{c|s|}\) and
\(\mathcal N_N^\sigma=O(\|\rho^\sigma\|^2)\), with the same three
solution-wise \(C^1\) parameter operators.
\end{lemma}

\begin{proof}
The attracting orbit is flushed through a fixed physical subannulus before
it reaches \(-S_\delta\).  The causal Green column in
\cref{eq:anchor-annulus-capped-action} therefore puts its normal defect at
\(-S_\delta\), together with its first parameter derivatives, below
\(C\mathfrak f_\delta\).

Subtract the exact invariance equation for \(H\) from the raw normal equation.
On \(J_{\sigma,S}\) the defect satisfies
\begin{equation}
 \mathcal L_{\rho,\sigma}\rho^\sigma
 =\mathcal N_{\rho,\sigma}(\rho^\sigma),\qquad
 \|D_\rho\mathcal N_{\rho,\sigma}\|
 \le C\|\rho^\sigma\|_{\mathcal X_{\sigma,S}^\sharp},          \label{eq:anchor-central-defect-equation}
\end{equation}
The generated-core normal Green operator is the
restriction of the physical current dichotomy to these fixed retained
intervals; its norm and its first strong-to-weak parameter derivative lose
at most a fixed power of \(\delta^{-1}\).  Apply its causal version first on
\(J_{a,S}\) to the flat input at \(-S_\delta\).  The contraction on the ball
\(\|\rho^a\|\le2C\mathfrak f_\delta\) gives the asserted bound through
\(s=0\).

The nonlinear chart identities
\cref{eq:anchor-central-chart-identities,eq:anchor-paired-stable-point}
now give the exact equality
\begin{equation}
 \mathfrak n_p(A^0)=\psi_A=\mathfrak n_p(B^0),
 \qquad \rho_0^a=\rho_0^r .                           \label{eq:anchor-paired-normal-coordinate}
\end{equation}
Thus the repelling defect problem has the already bounded input history at
zero.  Its forward normal Green problem gives the same bound on \(J_{r,S}\).
Differentiating the exact equality above and using the averaged derivative
of \cref{eq:anchor-central-defect-equation} proves the two first-parameter
bounds.  The conclusion concerns this paired point, not an arbitrary point
of the stable sheet.  This is a normal, solution-wise Green problem on the
exact physical graph; no ambient inverse is inferred from
\cref{lem:shared-graph-trace}, and no claim is made outside its uncut tube.

The central equations of the anchored and original models agree exactly.
Substituting the defining defect identity above into their raw planar
coordinate equation and applying Taylor's formula in the normal variable
gives \cref{eq:anchor-actual-field-expansion}.  The coefficient bounds
follow from the \(C^{12}\) raw chart and the retained-tube envelopes.
Accordingly, the perturbation used in every first-integral identity is
\[
 Q^{\sigma,\rm anc}
 =z^{\sigma,\rm anc\,\prime}-q_0(z^{\sigma,\rm anc})
 =\delta q_{1,N}+\delta^2q_{2,N}
  +\delta^3\mathcal R_{3,N}
  +\mathcal B_N^\sigma\rho^\sigma
  +\mathcal N_N^\sigma(\rho^\sigma).
\]
\end{proof}

\begin{lemma}[Anchored endpoint-tail completion]
\label[lemma]{lem:anchor-endpoint-tail-completion}
Put \(R=S_\delta+B\), and let \(z^{a,\rm anc}\) and \(z^{r,\rm anc}\)
be the planar core traces of the exact paired histories through
\(A^0\) and \(B^0\).  Define
\begin{align}
 \mathcal B_{\rm out}:={}&
 c_{\mathscr H}\{\mathscr H_\alpha(z^{a,\rm anc}(-R))
       -\mathscr H_\alpha(z^{r,\rm anc}(R))\}\notag\\
 &+c_{\mathscr H}\int_{-R}^{-S_\delta}
   \nabla\mathscr H_\alpha(z^{a,\rm anc})^T
       Q^{a,\rm anc}\,\dd s
 +c_{\mathscr H}\int_{S_\delta}^{R}
   \nabla\mathscr H_\alpha(z^{r,\rm anc})^T
       Q^{r,\rm anc}\,\dd s .                     \label{eq:anchor-outer-endpoint-strip}
\end{align}
Here \(Q^{\sigma,\rm anc}=z^{\sigma,\rm anc\,\prime}
-q_0(z^{\sigma,\rm anc})\) is the exact planar perturbation of the actual
RFDE trace.  Define the linear structural profile by
\[
 a^{\rm anc}_{2,N}(s)[\eta]
 =-e^{-s^2/2}s^2\Lambda_N(\eta).
\]
It is oriented so that
\[
 \int_{\mathbb R}a^{\rm anc}_{2,N}(s)[\eta]\,\dd s
 =-\sqrt{2\pi}\Lambda_N(\eta).
\]
Then
\begin{align}
 \mathcal B_{\rm out}
={}&\delta\int_{|s|>S_\delta}
       a^{\rm anc}_{1,N}(s,\nu)\,\dd s
 +\delta^2\int_{|s|>S_\delta}
       a^{\rm anc}_{2,N}(s)[\eta]\,\dd s\notag\\
 &+E_{\partial,N}^{\rm alg}+E_{\partial,N}^{\rm anc}.          \label{eq:anchor-endpoint-tail-completion}
\end{align}
For \(\mathfrak A_1=\{1,\partial_\nu,D_\eta\}\),
\begin{align}
 \sum_{\mathscr D\in\mathfrak A_1}
   \|\mathscr D E_{\partial,N}^{\rm anc}\|
 &\le C\mathfrak f_\delta,\notag\\
 |E_{\partial,N}^{\rm alg}|
 +|\partial_\nu E_{\partial,N}^{\rm alg}|
 &\le C\delta^2,\notag\\
 \|D_\eta E_{\partial,N}^{\rm alg}\|
 &\le C\{\delta^3+\delta^2\|\eta\|\}.             \label{eq:anchor-endpoint-tail-errors}
\end{align}
In particular, the physical endpoint values are Gaussian-tail carriers;
they are not individually \(O(\mathfrak f_\delta)\).
\end{lemma}

\begin{proof}
Let \(s_{\rm loc}^\pm\) be the fold-coordinate cuts at
\(r=\pm r_{\rm loc}\).  They satisfy
\(|s_{\rm loc}^\pm|\ge c/\delta\).  On the interval between these cuts the
anchored and original equations are identical.  The raw original-endpoint
chart fixes the \(s\)-clock and exact planar coordinate on
\([s_{\rm loc}^+,0]\) and \([0,s_{\rm loc}^-]\).  The annular mixed
inverse and the three smoothing buffers in
\cref{lem:anchor-annulus-Green-graphs} put both exact histories in this
strong chart with one strong-to-weak parameter derivative.

Write
\[
 \mathcal I^\sigma(s)
 =c_{\mathscr H}\nabla\mathscr H_\alpha(z^{\sigma,\rm anc}(s))^T
       Q^{\sigma,\rm anc}(s).
\]
Direct substitution in the raw physical chart gives the following
single-Gaussian ledger on
\(S_\delta<|s|<|s_{\rm loc}^\pm|\):
\begin{align}
 \mathcal I^\sigma(s)
  -\delta a^{\rm anc}_{1,N}(s,\nu)
  -\delta^2a^{\rm anc}_{2,N}(s)[\eta]
 &=e^{-s^2/2}\widehat{\mathcal E}_{\rm alg}^\sigma(s)
      +\mathcal E_{\rm anc}^\sigma(s),             \label{eq:anchor-outer-integrand-decomposition}\\
 |\widehat{\mathcal E}_{\rm alg}^\sigma(s)|
  +|\partial_\nu\widehat{\mathcal E}_{\rm alg}^\sigma(s)|
 &\le C\delta^2e^{c|s|}\langle s\rangle^M,\notag\\
 \|D_\eta\widehat{\mathcal E}_{\rm alg}^\sigma(s)\|
 &\le C\{\delta^3+\delta^2\|\eta\|\}
          e^{c|s|}\langle s\rangle^M,              \label{eq:anchor-outer-coefficient-ledger}\\
 \sum_{\mathscr D\in\mathfrak A_1}
  \int_{S_\delta<|s|<|s_{\rm loc}^\pm|}
       \|\mathscr D\mathcal E_{\rm anc}^\sigma(s)\|\,\dd s
 &\le C\mathfrak f_\delta .                        \label{eq:anchor-outer-boundary-ledger}
\end{align}
There is exactly one Gaussian factor in
\cref{eq:anchor-outer-integrand-decomposition}.  The hatted remainder is a
de-Gaussianized coefficient remainder, not the first-integral integrand
itself.

To verify the ledger, use the \(C^1\) actual-parameter original-endpoint
family and its physical-speed bordered inverse to bound the raw trace and
its first parameter derivative.  The coefficient allocation before any
integration is
\begin{align}
 c_{\mathscr H}\nabla\mathscr H_\alpha(\gamma_0)^T
       \delta q_{1,N}(\gamma_0,\nu)
   &=\delta a^{\rm anc}_{1,N}(s,\nu),\notag\\
 c_{\mathscr H}\nabla\mathscr H_\alpha(\gamma_0)^T
       \delta^2D_\eta q_{2,N}(\gamma_0,\nu,0)[\eta]
   &=\delta^2a^{\rm anc}_{2,N}(s)[\eta],            \label{eq:anchor-endpoint-coefficient-table}\\
 \mathcal E_{\rm tr}^\sigma
   &:=c_{\mathscr H}
     \{\nabla\mathscr H_\alpha(z^{\sigma,\rm anc})^T
          q_{1,N}(z^{\sigma,\rm anc},\nu)
       -\nabla\mathscr H_\alpha(\gamma_0)^T
          q_{1,N}(\gamma_0,\nu)\},\notag\\
 \mathcal E_{\rm str}^\sigma
   &:=c_{\mathscr H}\nabla\mathscr H_\alpha
        (z^{\sigma,\rm anc})^T
       \{q_{2,N}(z^{\sigma,\rm anc},\nu,\eta)
          -D_\eta q_{2,N}(\gamma_0,\nu,0)[\eta]\},\notag\\
 \mathcal E_{\rm jet}^\sigma
   &:=c_{\mathscr H}\nabla\mathscr H_\alpha
        (z^{\sigma,\rm anc})^T\delta^3\mathcal R_{3,N}.
\end{align}
Taylor's formula and the actual trace bounds give, with the Gaussian removed,
\begin{align}
 \sum_{\mathscr D\in\{1,\partial_\nu\}}
   \{|\mathscr D(\delta\mathcal E_{\rm tr}^\sigma)|
      +|\mathscr D(\delta^2\mathcal E_{\rm str}^\sigma)|\}
   &\le C\delta^2e^{c|s|}\langle s\rangle^M,\notag\\
 \|D_\eta(\delta\mathcal E_{\rm tr}^\sigma)\|
   +\|D_\eta(\delta^2\mathcal E_{\rm str}^\sigma)\|
   &\le C\{\delta^3+\delta^2\|\eta\|\}
          e^{c|s|}\langle s\rangle^M,\notag\\
 \sum_{\mathscr D\in\mathfrak A_1}
   |\mathscr D\mathcal E_{\rm jet}^\sigma|
   &\le C\delta^3e^{c|s|}\langle s\rangle^M.         \label{eq:anchor-endpoint-coefficient-errors}
\end{align}
The terms
\(\mathcal B_N^\sigma\rho^\sigma+
  \mathcal N_N^\sigma(\rho^\sigma)\) in
\cref{eq:anchor-actual-field-expansion} are assigned to
\(\mathcal E_{\rm anc}^\sigma\), not counted again in the hatted algebraic
remainder.  Their three \(C^1\) solution-wise integrals are bounded by
\cref{eq:anchor-actual-graph-defect-flat}.  The part depending on the anchor
collar is propagated through a fixed physical subannulus and has the same
bound.  This proves
\cref{eq:anchor-outer-coefficient-ledger,eq:anchor-outer-boundary-ledger}
term by term.  The calculation uses
\cref{eq:shared-q1,eq:q2-physical-continuation} only for the two displayed
physical profiles; it does not use
\cref{lem:anchor-actual-central-defect} outside its retained tube and does
not use a selected endpoint trace.

Apply
\[
 \frac{d}{ds}\mathscr H_\alpha(z(s))
 =\nabla\mathscr H_\alpha(z(s))^T
       \{z'(s)-q_0(z(s))\}
\]
first from \(s_{\rm loc}^+\) to \(-R\) on the attracting trace and then
from \(R\) to \(s_{\rm loc}^-\) on the repelling trace.  The first-integral
values at the two \(r_{\rm loc}\)-cuts split into the singular critical
value and the boundary-dependent normal column.  The first is
\(O(e^{-c/\delta^2})\); the second is bounded by
\cref{eq:anchor-annulus-capped-action}.  Together with
\cref{eq:anchor-outer-boundary-ledger}, these terms form
\(E_{\partial,N}^{\rm anc}\).

Add and subtract the explicit physical profiles on
\(S_\delta<|s|<|s_{\rm loc}^\pm|\).  Their portions beyond the two
\(r_{\rm loc}\)-cuts are \(O(e^{-c/\delta^2})\), with the three displayed
\(C^1\) parameter operators.
Consequently the endpoint values plus the two outer-strip integrals equal
the profile integrals over \(|s|>S_\delta\), with the signs in
\cref{eq:anchor-endpoint-tail-completion}.  Integrating the single Gaussian
in \cref{eq:anchor-outer-integrand-decomposition} and using
\begin{equation}
 \int_{S_\delta}^{\infty}
 e^{-s^2/2+cs}\langle s\rangle^M\,\dd s
 \le C e^{-S_\delta^2/2+cS_\delta}
       \langle S_\delta\rangle^M                   \label{eq:anchor-single-Gaussian-tail}
\end{equation}
gives \cref{eq:anchor-endpoint-tail-errors}.  Thus endpoint terms and outer
strips are neither discarded nor double counted, and no selected
preparation is used.  This proves the lemma.
\end{proof}

\begin{lemma}[Paired exact traces and anchored finite-gap bridge]
\label[lemma]{lem:anchor-fold-gap-conjugacy}
Before the scalar root is sought, the inward unstable branch has a unique
forward hit of \(\Sigma_0\), depending \(C^1\) on \((\nu,\eta)\), and the
future-asymptotic mixed problem
represents \(W^s(E_N^-)\) as the graph
\cref{eq:anchor-stable-sheet-graph} on the same section throughout the common
parameter box.  Pairing the incoming point with the stable-sheet point in
\cref{eq:anchor-paired-stable-point} gives two exact RFDE traces and the
first-integral gap \cref{eq:anchor-physical-gap}.  It has the zero-fiber
interpretation in \cref{prop:anchor-gap-comparison}, and its exact
finite-section decomposition gives
\cref{eq:anchor-physical-gap-estimates}
without a cokernel/domain identification and without an admissible
preparation.
\end{lemma}

\begin{proof}
The incoming orbit lies in the raw-compatible attracting chart by
\cref{thm:anchor-annulus-flat-forgetting}.  Eventual smoothing over the
reserved delay buffers and the strict physical speed sign give a \(C^1\)
parameter family of forward hits of \(X=0\).  The moving-time derivative is divided only by the
uniformly nonzero section speed.  No assertion that an off-root orbit already
traverses the repelling channel is used.

On the future side, the infinite mixed Green problem in
\cref{lem:anchor-annulus-Green-graphs} solves every stable entry coordinate
and removes the one unstable coefficient by future boundedness.  Its
intrinsic entry graph \cref{eq:anchor-intrinsic-entry-sheet}, pulled back by
the forward membership composition
\cref{eq:anchor-forward-preimage-membership}, gives
\(u=F^{-,0}(\psi)\) before a gap is evaluated.  The direct column identity
\cref{eq:anchor-forward-column-coefficient-identity} proves the required
transversality.  This is the explicit ambient range construction used here;
the graph-tangent operator in \cref{lem:shared-graph-trace} is not misused as
an ambient RFDE inverse.

By \cref{eq:anchor-central-chart-identities,eq:anchor-paired-stable-point},
the two points \(A^0,B^0\) have the same exact normal/history coordinate and
their normalized first-integral coordinates are, respectively,
\(u_A\) and \(F^{-,0}(\psi_A)\).  Consequently
\begin{equation}
 G^{\rm anc}=u_A-F^{-,0}(\psi_A)                     \label{eq:anchor-gap-exact-coordinate-identity}
\end{equation}
with no Taylor multiplier.  Since \(\mathfrak C_{0,p}\) is injective,
\(G^{\rm anc}=0\) if and only if \(A^0=B^0\), which is exactly membership
of the incoming history in the future stable sheet.  Its past is the
unstable orbit from \(E_N^+\), and its future is supplied by the stable
sheet.  RFDE forward uniqueness gives one global heteroclinic.  Conversely
every heteroclinic using the inward branch and lying in
\(\mathcal U_{\rm anc}\) has this fixed-phase hit and therefore has zero
gap.  This proves only the declared local uniqueness, not uniqueness among
all global connections.

Let \(z^{a,\rm anc}\) and \(z^{r,\rm anc}\) be the planar coordinates of
the exact orbit segments through \(A^0\) and \(B^0\), and define
\[
 Q^{\sigma,\rm anc}(s)
 =\frac{d}{ds}z^{\sigma,\rm anc}(s)-q_0(z^{\sigma,\rm anc}(s)).
\]
The definition \cref{eq:anchor-physical-gap} and the fundamental theorem of
calculus give the exact identity
\begin{align}
 G^{\rm anc}={}&c_{\mathscr H}\{\mathscr H_\alpha
       (z^{a,\rm anc}(-R))
       -\mathscr H_\alpha(z^{r,\rm anc}(R))\}\notag\\
 &+c_{\mathscr H}\int_{-R}^0\nabla\mathscr H_\alpha
       (z^{a,\rm anc}(s))^TQ^{a,\rm anc}(s)\,\dd s\notag\\
 &+c_{\mathscr H}\int_0^R\nabla\mathscr H_\alpha
       (z^{r,\rm anc}(s))^TQ^{r,\rm anc}(s)\,\dd s
                                                        \label{eq:anchor-exact-finite-gap-identity}
\end{align}
with no additional ambient term: all transverse histories enter the exact
fields \(Q^{\sigma,\rm anc}\).  Splitting the two integrals at
\(\pm S_\delta\) gives, literally,
\begin{align}
 G^{\rm anc}
 =\mathcal B_{\rm out}
 &+c_{\mathscr H}\int_{-S_\delta}^{0}
   \nabla\mathscr H_\alpha(z^{a,\rm anc})^T
      Q^{a,\rm anc}\,\dd s\notag\\
 &+c_{\mathscr H}\int_{0}^{S_\delta}
   \nabla\mathscr H_\alpha(z^{r,\rm anc})^T
      Q^{r,\rm anc}\,\dd s .                     \label{eq:anchor-exact-central-outer-split}
\end{align}
This equation both fixes all signs and shows that the outer strips are
counted once.

The coefficient transfer on these actual, not prepared, traces is
\cref{lem:anchor-actual-central-defect}.  It is solution-wise after three
smoothing buffers and does not assert operator differentiability on rough
histories.  The planar one-sided trace inverses in
\cref{lem:shared-graph-trace} are used only for the planar equations with
the last two terms as known sources.  Their Gaussian pairing norm and
\cref{eq:anchor-actual-graph-defect-flat} give the same trace estimates as
\cref{eq:shared-trace-jets} for the value and first parameter derivatives,
up to \(C\mathfrak f_\delta\).

On \(|s|\le S_\delta\), the anchored and original physical equations agree
literally.  The coefficient calculation
\cref{eq:shared-leading-gap,eq:shared-melnikov} therefore applies to the
integrands in \cref{eq:anchor-exact-finite-gap-identity}, giving
\begin{align}
 \delta\sqrt{2\pi}(\nu-\nu_{0,N}),\qquad
 -\delta^2\sqrt{2\pi}\Lambda_N(\eta)               \label{eq:anchor-physical-Melnikov-pairings}
\end{align}
for the first parameter and structural pairings.  The half-line Green
estimates put only the anchor-dependent part below
\(\mathfrak f_\delta\).  The endpoint values themselves are combined with
the two outer strips by
\cref{lem:anchor-endpoint-tail-completion}; this supplies exactly the
missing whole-line tails in
\cref{eq:anchor-physical-Melnikov-pairings}.  The residual Gaussian
tail error is bounded by
\begin{equation}
 C\delta^{-M}e^{-S_\delta^2/2+cS_\delta}
       \langle S_\delta\rangle^M=\mathfrak t_\delta.           \label{eq:anchor-physical-tail-bound}
\end{equation}
The moving-trace, parameter-Taylor, and graph remainders are estimated by
applying Taylor's formula to
\cref{eq:anchor-actual-field-expansion}.  The resulting integrals are the
same local multilinear expressions as
\cref{eq:gap-error-moving-nu,eq:gap-error-moving-zeta,%
eq:gap-error-jet-bound,eq:gap-error-parameter-derivatives}, restricted to
their value and first parameter derivatives, plus the
explicit flat defect in \cref{eq:anchor-actual-graph-defect-flat}.  Thus no
selected endpoint or selected trace estimate is transferred by assertion.
Enlarge \(M_{\rm anc}\), if necessary, to dominate the finite \(C^1\)
loss exponent \(M_*\).  Choosing \(p>M_{\rm anc}+6\) then absorbs
\(\mathfrak e_\delta=\mathfrak f_\delta+\mathfrak t_\delta\) below every
displayed algebraic order and yields
\cref{eq:anchor-physical-gap-estimates}.
This proves the lemma.
\end{proof}

\begin{proof}[Proof of \cref{prop:anchor-gap-comparison}]
Both the zero-fiber statement and the direct physical estimates are
\cref{lem:anchor-fold-gap-conjugacy}.
\end{proof}

\begin{proof}[Proof of \cref{thm:anchored-physical-root-conormal}]
Choose \(c_0\) strictly inside the common \(\nu\)-window.  The first two
lines of \cref{eq:anchor-physical-gap-estimates} give opposite signs at
\(\nu_{0,N}\pm c_0\) and
\begin{equation}
 \partial_\nu G^{\rm anc}
 \ge\frac12\delta\sqrt{2\pi}>0                     \label{eq:anchor-physical-parameter-transversality}
\end{equation}
after decreasing \(\delta_0,\eta_0\).  The intermediate-value theorem and
strict monotonicity give the unique root.  By
\cref{prop:anchor-gap-comparison}, it is precisely the heteroclinic root of
the fixed physical equation, so its definition contains no preparation.
  On the attracting side the root orbit is the recut incoming curve
  \(\mathcal Q^+\).  On the repelling side, the entry graph radius
  \cref{eq:anchor-half-line-radius}, the fixed-point difference formula
  \cref{eq:anchor-tail-value-difference}, and forward event alignment on a
  fixed subannulus bound its distance from the specified chart-center curve
  \(\mathcal Q^-\) by
  \(CR_{\mathrm{hl},\delta}\le C\delta^2\) in the fading norm.  This
  proves the stated outer history bound against the two curves defined in the
  theorem, rather than against an unspecified slow graph.  Crossing a fixed \(r\)-subannulus under
\(\dot r=\delta^2q_0+O(\delta^4)\) takes at least
\(c\delta^{-2}\) original time.  On the central window the range inverse
for the planar trace, the actual normal-defect estimate
\cref{lem:anchor-actual-central-defect}, and
\cref{eq:shared-trace-jets} give
\cref{eq:anchor-heteroclinic-canard-tracking}.  At the root the two
one-sided histories are restrictions of the same orbit.

Implicit differentiation gives
\begin{equation}
 D_\eta\nu_c^{\rm anc}
 =-\frac{D_\eta G^{\rm anc}}
          {\partial_\nu G^{\rm anc}}
 =\delta\Lambda_N+O(\delta^2+\delta\|\eta\|).       \label{eq:anchor-root-implicit-derivative}
\end{equation}
Multiplication by \(\delta^2\) proves
\cref{eq:anchor-physical-root-jets}.  Integrating the first derivative
on the structural segment proves \cref{eq:anchor-physical-root-increment}.
Dividing by \(\delta^3\) gives
\cref{eq:anchor-physical-conormal-limit}.

Finally, \(\Lambda_N\) is already identified with the hidden-return Schur
factorization in \cref{eq:joint-schur-hyperplane}, and all the range,
tomography, one-delay, and witness statements concern that same covector.
They therefore transfer verbatim to the intrinsic limiting conormal.
\end{proof}

\begin{proof}[Proof of \cref{prop:anchor-physical-naturality-composition}]
A phase or section change is used only on the constructed incoming orbit and
future-asymptotic sheet.  Let \(g_1\) and \(g_2\) be the resulting scalar
parameter-space membership functions at two regular cuts.  Forward uniqueness
and event recutting give
\[
 g_1(p)=0\quad\Longleftrightarrow\quad g_2(p)=0
\]
on the common local parameter box.  Both differentials are nonzero at the
root.  Use \(g_1\) as one coordinate in a Banach product chart.  In those
coordinates the fundamental theorem of calculus in the scalar coordinate
gives \(g_2=a g_1\) with \(a\) continuous and
\(a(p_c)\ne0\).  Hence \(Dg_2(p_c)=a(p_c)Dg_1(p_c)\), so the conormal line is
unchanged.  This argument uses no inverse history holonomy.  Multiplication by
a nowhere-zero defining factor is the same calculation, and a nonsingular
parameter change acts by the ordinary covector pullback.

For intermediate sections write the total forward holonomy on the constructed
solution families as \(P=P_j\circ\cdots\circ P_1\).  The forward chain rule
gives \(DP=DP_j\cdots DP_1\).  Each cut produces the same local parameter
zero locus, so the preceding defining-function argument makes its conormal
line independent of inserting or deleting a cut.  This is the
complete-history transition-jet composition law; no factor \(DP_i^{-1}\) is
formed.

For each anchor multiplier, apply the direct gap estimates
\cref{eq:anchor-physical-gap-estimates} and the implicit derivative
\cref{eq:anchor-root-implicit-derivative}.  Their constants are uniform over
the bounded completion class because the clock lower envelope is part of
\cref{eq:anchor-admissible-class}.  After centering at its own baseline,
each weighted conormal is therefore within \(C\delta\) of
\((1,-\Lambda_N)\).  The triangle inequality gives the stated
\(2C\delta\) pairwise bound.  This argument does not claim a pairwise flat
gap comparison or equality of exact roots.
\end{proof}
